% ====================================================================
% Chapter 1: Graphite anode ICA tail theory
% Scope: discharge/delithiation branch; thermodynamic stage + kinetic tail.
% ====================================================================
\documentclass[11pt,a4paper]{article}
\usepackage[margin=22mm]{geometry}
\usepackage{kotex}
\usepackage{amsmath,amssymb,mathtools,bm}
\usepackage{booktabs,longtable,array}
\usepackage{enumitem}
\usepackage{amsthm}
\usepackage{hyperref}
\usepackage{fancyhdr}
\usepackage{lastpage}
\usepackage{setspace}
\setstretch{1.13}
\setlength{\parskip}{0.45em}\setlength{\parindent}{0pt}\setlength{\headheight}{14pt}
\hypersetup{colorlinks=true, linkcolor=blue, urlcolor=blue, citecolor=blue,
  pdftitle={Graphite anode ICA tail theory - Chapter 1}}
\pagestyle{fancy}\fancyhf{}
\lhead{Graphite anode ICA tail theory --- Ch.1}\rhead{\thepage/\pageref{LastPage}}
\renewcommand{\headrulewidth}{0.4pt}

\newtheorem*{groundingbox}{Grounding}
\newtheorem*{boundbox}{유효범위(Validity bound)}
\newtheorem*{flagbox}{FLAGGED (미확정)}
\newtheorem*{keybox}{Keystone}

\newcommand{\dd}{\mathrm{d}}
\newcommand{\eff}{\mathrm{eff}}
\newcommand{\eq}{\mathrm{eq}}
\newcommand{\app}{\mathrm{app}}
\newcommand{\drive}{\mathrm{drive}}
\newcommand{\bg}{\mathrm{bg}}
\newcommand{\cell}{\mathrm{cell}}
\newcommand{\ext}{\mathrm{ext}}
\newcommand{\OCV}{\mathrm{OCV}}
\newcommand{\tot}{\mathrm{tot}}
\newcommand{\AL}[1]{\textsf{[AL-#1]}}

\title{\textbf{리튬이온전지 흑연 음극 ICA($dQ/dV$) tail 이론 --- Chapter 1}\\[0.4em]
\large Thermodynamic stage, electrode-potential-assisted relaxation barrier, and first-pass fitting approximation}
\author{}\date{}

\begin{document}
\maketitle

% ====================================================================
\section*{서 — 하나의 관찰에서 출발하는 단일 인과 사슬}
\addcontentsline{toc}{section}{서 — 단일 인과 사슬}
% ====================================================================
본 장은 다음 \textbf{하나의 관찰}에서 출발하여 끝까지 그 끈을 놓지 않는다.
\begin{quote}\itshape
흑연 음극의 $dQ/dV$(ICA, incremental capacity) 곡선에서 상변이 peak 의 꼬리(tail)가 \textbf{온도가
낮을수록 길게 늘어져 다음 peak 와 겹치고}, 높을수록 짧게 끝난다. 또 전류(C-rate)가 커지면 peak 겹침이
심해진다.
\end{quote}
이 관찰을 \textbf{다섯 단계}로 푼다. 각 단계는 학부 수준(화학열역학·전기화학·미적분)으로 따라갈 수 있게
중간을 생략하지 않는다.
\begin{enumerate}[label=\textbf{(\arabic*)},leftmargin=2.2em]
\item \textbf{평형 열역학만으로는 설명이 안 된다.} 뒤에서 보이듯 평형 전이폭은 $w=RT/F$ 라 \emph{저온서
  오히려 좁아진다}(\S\ref{sec:eq}). 평형 이론은 ``저온 $\to$ 짧은 꼬리''를 예측한다 — 관찰과 \emph{반대}.
\item \textbf{따라서 꼬리의 주역은 동역학이다.} 진행률 $\xi_j$ 가 평형 목표 $\xi_{\eq,j}$ 를 \emph{유한
  속도}로 좇으며 생기는 지연(lag)이 꼬리를 만든다. 속도상수 $k\propto e^{-\Delta G_a/RT}$ 가 저온서
  작아져 지연이 커지고 꼬리가 길어진다(\S\ref{sec:dyn}) — 관찰과 \emph{일치}.
\item \textbf{그 배리어는 극판 전위가 낮춘다.} 전류가 흐를 때 내부 전위는 평형(OCV)과 다르며, 그 전위
  구동력 $\mathcal A_j$ 가 유효 배리어를 낮춰 진행 속도를 바꾼다(\S\ref{sec:rate}). 그래서 C-rate 변화는 potential-assisted rate 변화와 apparent-axis polarization을 통해 peak overlap을 바꾼다.
\item \textbf{무대는 전하 보존이 깐다.} 내부 음극 전위 $V_n$ 은 외부에서 읽어오는 값이 아니라 \emph{전하
  보존식의 해}로 결정된다(\S\ref{sec:charge}). 꼬리의 \emph{깊이}는 배리어 분포가 만드는
  \emph{relaxation-length spectrum} 의 중첩(kernel integral)으로 설명된다(\S\ref{sec:spectrum}--\ref{sec:kernel}).
\item \textbf{결과 = 피팅 가능한 닫힌 논리식}(\S\ref{sec:fiteq}). 이후 발열(Ch.2,4)$\cdot$반응속도론
  (Ch.3)$\cdot$히스테리시스(Ch.5)로 확장하고, 수치$\cdot$피팅 실무는 \emph{본 장 안}(\S\ref{sec:ch1_simplefit} 심플
근사식, \S\ref{sec:planB}--\ref{sec:planA} closure, 부록 \S\ref{sec:ch1_numeric})에서 닫는다.
\end{enumerate}
\noindent\textbf{한 문장 요지: 열역학이 \emph{무대}($V_n$ 과 평형 target)를 깔고, 관찰된 꼬리 거동의
\emph{주역}은 동역학이다.} 본 Chapter 1 은 \textbf{방전(탈리튬화) 방향}만 다룬다. 충전 branch,
히스테리시스, full-cell 전압 분해는 Ch.5 에서 정의한다.

\begin{groundingbox}
\textbf{지배 표준(GS).} \textbf{(GS-1)} 모든 식은 물리적 의미를 먼저 말로 제시한 뒤 수식화하고 중간
단계를 생략하지 않는다(학부 무비약). \textbf{(GS-2)} 결론은 출판 수준 엄밀성. \textbf{(GS-3)} 모든 가정은
통합 Assumption Ledger(\AL{\#}: 논문/교과서 근거 + 검증 DOI)를 인용하며, 근거가 없으면 FLAGGED 로
표기하고 established 로 쓰지 않는다. \textbf{(GS-4)} 임의 clip/cap/softplus/step-function 정의식과 근거
없는 수치 임계값을 물리 모델로 쓰지 않는다(속도 제한은 물리적 병목으로 도입). 수치 해법과 피팅 실무는
본 장 부록 \S\ref{sec:ch1_numeric} 로 분리한다.
\end{groundingbox}

\tableofcontents
\newpage

% ====================================================================
\section{기호와 단위}\label{sec:notation}
% ====================================================================
진행률 기호는 $\xi$ 로 통일한다(일부 문헌의 $\theta$ 는 같은 occupation/progress variable로만 언급한다). 방전 누적 전하 좌표 $q$, 외부 전류 $I$,
방향 부호 $s_I,s_{\phi,j}\in\{+1,-1\}$(방전 기본 $+1$).
\begin{longtable}{p{0.16\textwidth} p{0.10\textwidth} p{0.66\textwidth}}
\toprule 기호 & 단위 & 의미 \\ \midrule \endhead
$q$ & --- & 방전 진행 좌표, $q=Q_\ext/Q_\cell$ \\
$Q_\ext,\ Q_\cell$ & C & 외부 누적 방전 전하, 기준 셀 용량(쿨롱) \\
$V_n$ & V & 내부 음극 전위 — 전하 보존식의 해(\AL{1}) \\
$V_{n,\app}$ & V & 측정되는 apparent 전위 $=V_n+s_I|I|R_n$ (\AL{6}) \\
$V_{n,\drive}$ & V & 속도식 구동 전위; 기본 $=V_n$ \\
$V_{n,\OCV}$ & V & 평형($\xi_j=\xi_{\eq,j}$)에서의 전하보존식 해 \\
$R_n$ & $\Omega$ & 음극 유효 분극 계수(전압축 이동) \\
$\xi_j,\ \xi_{\eq,j}$ & --- & $j$번째 effective transition 진행률, 평형 진행률 ($0\!\le\!\xi\!\le\!1$) \\
$U_j(T),\ w_j(T)$ & V & 전이 중심 전위(\S\ref{sec:eq} 에서 $U_j\equiv U_j^0-\mu^0/(s_{\phi,j}F)$ 로 $\mu^0$ 흡수), 평형 전이 폭 \\
$Q_{j,\tot}$ & C & $j$번째 transition 의 방전 방향 용량 기여 ($>0$) \\
$Q_\bg(V_n,T)$ & C & staging 으로 분리 안 되는 배경 누적 용량; $\partial Q_\bg/\partial V_n$=잔류 chemical capacitance \\
$\mathcal A_j$ & J/mol & 상변이 구동력 $=s_{\phi,j}F(V_{n,\drive}-U_j)$ \\
$\chi_j$ & --- & Level-A scalar relaxation-barrier sensitivity ($0\le\chi_j\le1$). Ch.1 에서는 Ch.3 의 transfer coefficient $\beta_j$ 와 동일시하지 않음 \\
$\Delta G_{a,j},\Delta H_{a,j},\Delta S_{a,j}$ & J/mol & 활성화 자유에너지/엔탈피/엔트로피 \\
$G$ & J/mol & local activation barrier (mode 별 값; \S\ref{sec:spectrum} 에서 분포) \\
$k_j,\ k_0(T)$ & 1/s & 진행률 완화 속도상수(mobility), Eyring prefactor \\
$L$ & --- & local relaxation length $=|I|/(Q_\cell k_j)$ \\
$\rho_G(G;T)$ & mol/J & activation barrier 확률밀도 ($\int_0^\infty\rho_G\,\dd G=1$; $G$ 가 J/mol 이므로 단위는 그 역수 mol/J) \\
$A_L^{\mathrm{prob}},A_L^{\mathrm{amp}}$ & --- & relaxation-length 확률밀도(정규화) / 진폭 스펙트럼 $=A_0 A_L^{\mathrm{prob}}$; kernel 의 $A_L\equiv A_L^{\mathrm{amp}}$ (\S\ref{sec:spectrum}) \\
$A_0(L)$ & --- & mode 별 물리적 진폭 가중(입자크기분포$\cdot$활성 site 밀도 등 독립 물리량) \\
$L_0,\ L_{\min}$ & --- & $L_0=|I|/(Q_\cell k_0)$ (prefactor 기준 길이), $L_{\min}=L_0 e^{-\chi_j\mathcal A_j/RT}$ (support 하한, $G\!\ge\!0$) \\
$L_\varphi$ & V & 전압축 특성 꼬리 길이 $=|\dd V_n/\dd q|_{q_a}L$ (\S\ref{sec:fiteq}, 식~\eqref{eq:Lphi}) \\
$\Theta,\ Q_p$ & ---,\ C & 전체 effective phase progress, 그 용량 scale \\
\bottomrule
\end{longtable}
\textbf{단위 주의.} 미분방정식의 전류 $I$ 가 A $=\mathrm{C/s}$ 이므로 $Q_\cell$ 은 쿨롱이어야 한다
($Q_\cell[\mathrm C]=3600\,Q_\cell[\mathrm{Ah}]$). 외부 전하·진행 좌표는 $Q_\ext(t)=\int_0^t|I|\,\dd t'$,
$q=Q_\ext/Q_\cell$, 정전류에서 $\dd q/\dd t=|I|/Q_\cell$.

% ====================================================================
\section{Convention과 assumption ledger}\label{sec:convention}
% ====================================================================
본 장은 문헌의 graphite lithiation stage order를 그대로 fitting label로 쓰지 않는다. Graphite literature에서는
stage evolution을 lithiation order로 기술하는 경우가 많지만, 여기서는 measured discharge/delithiation branch에서 관측되는
peak와 tail을 해석한다. 따라서 $j$는 crystallographic stage number가 아니라 \textbf{branch-local effective transition label}이다.
각 $j$에서 $\xi_j=0\to1$은 해당 discharge branch의 local transition progress가 진행되는 방향을 뜻한다.

\textbf{Barrier convention.} $G$ 또는 $\Delta G_{a,j}$는 local mode의 intrinsic relaxation barrier다.
$\Delta G_{\eff,j}=G-\chi_j\mathcal A_j$는 Ch.1에서 scalar relaxation mobility를 바꾸는 Level-A effective barrier로만 쓴다.
이는 microscopic forward barrier lowering 또는 Butler--Volmer transfer coefficient와 동일한 말이 아니다. Forward/backward 방향성을
가진 rate splitting은 Ch.3에서 별도 Level-B extension으로 도입한다.

\textbf{ICA fitting convention.} 본 장의 first-pass fitting은 measured ICA가 정확히 exponential이라는 주장이 아니다.
먼저 charge conservation으로 internal potential $V_n$과 apparent axis $V_{n,\app}$를 분리하고, small-tail/local-baseline window에서
baseline-subtracted signal만 $Y_{\mathrm{tail}}\simeq B\exp[-(q-q_a)/L]$로 근사한다. 이때 $B$는 scale/nuisance이고, shape information은
주로 $L$에 들어간다.

\begin{longtable}{p{0.10\textwidth} p{0.18\textwidth} p{0.30\textwidth} p{0.32\textwidth}}
\toprule ID & status & source type & allowed use \\ \midrule \endhead
\AL{1} & established & porous-electrode charge conservation \cite{doyle1993,newman} & $V_n$은 external lookup이 아니라 charge balance의 해다. \\
\AL{2} & established/bounded & lattice-gas and regular-solution thermodynamics \cite{mckinnon1983,bazant2013} & logistic target은 ideal limit, effective width는 calibrated coarse-graining. \\
\AL{3} & established & graphite staging observations \cite{dahn1991,ohzuku1993,funabiki1999ea,funabiki1999jes} & staging/effective transition이 ICA peak를 만든다. \\
\AL{4} & bounded calibration & low-rate OCV/GITT & $U_j,w_j,Q_{j,\tot},Q_{\bg}$는 equilibrium-stage fit에서 먼저 제한한다. \\
\AL{5} & validity gate & positive differential capacity & local fitting window에서 $C_{\bg}>0$ 및 positive ICA denominator가 필요하다. \\
\AL{6} & established & ohmic/transport polarization separation \cite{newman,bardfaulkner} & $V_{n,\app}$와 $V_n$를 분리해야 $\chi_j$ 해석이 가능하다. \\
\AL{7} & mathematical gate & range and monotonicity condition & charge-balance solution은 admissible window에서만 존재한다. \\
\AL{10} & established & transition-state/Eyring theory \cite{eyring1935} & rate scale은 $\exp[-\Delta G^\ddagger/(RT)]$를 따른다. \\
\AL{11} & bounded & BEP/BV/Marcus small-drive analogy \cite{evanspolanyi1938,bardfaulkner,marcus1956} & Ch.1에서는 scalar relaxation barrier의 first-order potential assistance로만 사용한다. \\
\AL{12} & established/bounded & Onsager/two-state local relaxation \cite{onsager1931,degrootmazur1962} & local first-order relaxation toward a target. \\
\AL{13} & flagged/bounded & heterogeneous electrode plus cooperative staging literature & barrier-only independent mode distribution is a reduced mean-field ansatz. \\
\AL{14} & established mechanism, bounded application & relaxation-time distributions \cite{svare2000,johnston2006,lindsey1980} & barrier distribution can create broad tails; graphite ICA application is model contribution. \\
\AL{15} & established bound & Marcus curvature \cite{marcus1956} & linear barrier assistance is a small-drive approximation. \\
\AL{16} & bounded analogy & Fredholm ratio-substitution methods \cite{lee2011,son2013,jcp2017} & unknown-ratio closure idea only; Fredholm formula is not copied into the Volterra problem. \\
\bottomrule
\end{longtable}
% ====================================================================
\section{흑연 staging 과 effective transition}\label{sec:staging}
% ====================================================================
흑연의 lithiation literature에서는 stage $4\to3\to2\mathrm L\to2\to1$ 같은 order가 자주 쓰이며, 각 stage transformation이 $dQ/dV_n$에 peak를
남긴다 \cite{dahn1991,ohzuku1993,funabiki1999ea,funabiki1999jes}(\AL{3}). 그러나 본 장의 fitting label은 measured discharge/delithiation branch 기준이다.
인접 전이의 중심 전위 차가 전이폭 이하이면 peak가 융합되어 관측 peak 수가 $N_p\approx3\sim4$가 되므로, 본 장은 분리 가능한 effective transition을
$j=1,\dots,N_p$로 둔다. 이 $j$는 crystallographic stage number가 아니라 branch-local peak label이다. 각 유효 transition 내부에도 입자 크기$\cdot$국소 환경$\cdot$결정립계
차이에 의한 \emph{local mode 분포}가 존재하며, 이것이 \S\ref{sec:spectrum} spectrum 의 물리적 근원이다.
진행률 방향은 방전 좌표 $q$ 와 같게 잡아 $\xi_j=0$(미진행)$\to\xi_j=1$(완료), $Q_{j,\tot}>0$.

% ====================================================================
\section{무대 (I): 전하 보존식과 내부 전위 $V_n$}\label{sec:charge}
% ====================================================================
\textbf{물리.} 외부 회로로 흘린 방전 전하와, 음극 내부의 방전 방향 누적 전하(배경 저장분 + staging 진행분)는
같아야 한다 — 단순한 전하 보존이다. 이 등식이 내부 전위 $V_n$ 을 결정한다(외부에서 주어지는 함수가
아니다; \AL{1}, \cite{doyle1993,newman}):
\begin{equation}
\boxed{\;Q_\cell\,q \;=\; Q_\bg(V_n,T) \;+\; \sum_{j=1}^{N_p} Q_{j,\tot}\,\xi_j\;.}
\label{eq:charge_balance}
\end{equation}
좌변은 외부 입력(아는 값), 우변은 내부 상태($V_n$ 과 $\xi_j$ 의 함수)다. 주어진 $(q,T,\{\xi_j\})$ 에서 이
식을 만족하는 $V_n$ 이 내부 전위다. 이는 다공성 전극 이론의 implicit 화학퍼텐셜 관계와 같은 구조다
\cite{doyle1993,newman}.

\textbf{배경 용량의 의미.} $Q_\bg(V_n,T)$ 는 단순 그래프 기저선이 아니라 뚜렷한 staging 으로 분리하지 않은
방전 방향 잔류 누적 용량이며, $\xi_j$ 가 평형 진행률을 즉시 따라가지 못할 때 전하 보존을 맞추기 위해 $V_n$
이 얼마나 이동해야 하는지를 정한다. 그 전위 기울기 $\partial Q_\bg/\partial V_n$ 이 잔류 chemical
capacitance 다. $V_n$ 을 안정적으로 풀려면 이 기울기가 양이어야 한다(\AL{5}; bazant2013 의 비평형
열역학적 미분용량):
\begin{equation}
\frac{\partial Q_\bg}{\partial V_n}\ \ge\ \epsilon_Q\ >\ 0 .
\label{eq:monotone}
\end{equation}
$\epsilon_Q$ 는 임의의 정규화 상수가 아니라 양의 미분용량이라는 물리 조건이며, 수치적으로는 해의 특이점을
막는 하한이다(\AL{5} BOUNDED).

\textbf{평형 극한 = OCV.} 평형(또는 충분히 느린 조건)에서는 $\xi_j=\xi_{\eq,j}(V_n,T)$ 이고
식~\eqref{eq:charge_balance} 가 $Q_\cell q=Q_\bg+\sum_j Q_{j,\tot}\xi_{\eq,j}$ 가 되어 그 해가 평형 음극
전위 $V_{n,\OCV}(q,T)$ 다. \textbf{즉 OCV 곡선은 임의로 주어진 lookup 이 아니라 전하 보존식의 평형
특수해다} — 중심식이 ``OCV 를 읽는'' 흐름으로 회귀하지 않는다.

\textbf{해 존재 조건(well-posedness).} 식~\eqref{eq:charge_balance} 가 $V_n$ 해를 가지려면 우변 목표가 배경
용량의 치역 안에 있어야 한다(\AL{7}): $Q_\cell q-\sum_j Q_{j,\tot}\xi_j\in\mathrm{Range}[Q_\bg(\cdot,T)]$.
단조성(식~\eqref{eq:monotone})과 이 치역 조건을 함께 만족하면 중간값 정리로 해가 유일하게 존재한다(부분
구간 피팅용 상대 기준형$\cdot$총용량 정합 등 실무는 \S\ref{sec:ch1_simplefit}).

% ====================================================================
\section{무대 (II): 평형 진행률 $\xi_\eq$ (logistic)}\label{sec:eq}
% ====================================================================
\textbf{왜 logistic 인가 (무비약 유도).} 흑연 intercalation 의 한 전이를 격자기체(lattice-gas)/정규용액
(regular-solution) 으로 보면, 점유율 $\xi$ 인 상태의 화학퍼텐셜은 \cite{mckinnon1983,bazant2013}(\AL{2})
\begin{equation}
\mu(\xi)=\mu^0+RT\ln\frac{\xi}{1-\xi}+\Omega_j(1-2\xi),
\label{eq:mu_latticegas}
\end{equation}
두 항의 출처를 학부 통계역학으로 한 줄씩 확인한다(식~\eqref{eq:mu_latticegas} 를 인용으로만 받지 않고 본문에서 유도한다).

\textbf{(i) 혼합 엔트로피 항.} $N$ 개의 동등한 자리 중 $n$ 개가 점유된 배열의 가짓수는
$W=\dfrac{N!}{n!\,(N-n)!}$ 이며, 점유율 $\xi\equiv n/N$ 으로 둔다. Boltzmann 관계 $S=k_B\ln W$ 에 Stirling
근사 $\ln m!\simeq m\ln m-m$ 를 적용하면(세 계승의 $-m$ 선형항이 $N=n+(N-n)$ 으로 상쇄되어)
\begin{equation}
S_{\mathrm{mix}}=k_B\big[N\ln N-n\ln n-(N-n)\ln(N-n)\big]
=-k_B N\big[\xi\ln\xi+(1-\xi)\ln(1-\xi)\big],
\label{eq:smix}
\end{equation}
즉 몰 기준($N\!\to\!N_A$, $R=k_B N_A$)으로 $S_{\mathrm{mix}}=-R[\xi\ln\xi+(1-\xi)\ln(1-\xi)]$ 다. 자유
에너지의 혼합 기여 $G_{\mathrm{mix}}=-TS_{\mathrm{mix}}=RT[\xi\ln\xi+(1-\xi)\ln(1-\xi)]$ 를 점유율로 미분하면
\begin{equation}
\frac{\partial G_{\mathrm{mix}}}{\partial\xi}=RT\big[(\ln\xi+1)-(\ln(1-\xi)+1)\big]=RT\ln\frac{\xi}{1-\xi}
\label{eq:mu_mix}
\end{equation}
($+1$ 과 $-1$ 이 상쇄). 여기서 $G_{\mathrm{mix}}$ 는 점유 자리 1몰당(intensive) 자유에너지이고, 화학퍼텐셜의 정의는
$\mu_{\mathrm{config}}=\partial G^{\mathrm{tot}}/\partial n$ ($G^{\mathrm{tot}}=N G_{\mathrm{mix}}$, 자리수 $N$, $\xi=n/N$)
인데 $\partial/\partial n=(1/N)\,\partial/\partial\xi$ 라 $N$ 인자가 상쇄되어 $\mu_{\mathrm{config}}=\partial G_{\mathrm{mix}}/\partial\xi$
다. 따라서 위 $\xi$-미분이 곧 화학퍼텐셜의 혼합 기여이며, 식~\eqref{eq:mu_latticegas} 의 첫 로그 항이 정확히 재현된다.

\textbf{(ii) 상호작용 항.} 정규용액 근사에서 점유 입자 간 상호작용(혼합 엔탈피)은
$G_{\mathrm{int}}=\Omega_j\,\xi(1-\xi)$ 이고, 그 미분 $\partial G_{\mathrm{int}}/\partial\xi
=\Omega_j(1-2\xi)$ (곱미분 $\dd[\xi-\xi^2]/\dd\xi=1-2\xi$)가 둘째 항이다. 표준 상태 항 $\mu^0$ 까지 더하면
$\mu=\mu^0+RT\ln[\xi/(1-\xi)]+\Omega_j(1-2\xi)$ 로 식~\eqref{eq:mu_latticegas} 가 완성된다(\AL{2};
$\mu^0,\Omega_j$ 의 미시값은 \cite{mckinnon1983,bazant2013} 에 위임).

\textbf{(iii) 전기화학 평형.} 평형은 이 화학퍼텐셜이 전극 전위와 균형을 이루는 조건
$\mu(\xi_\eq)=s_{\phi,j}F(V_n-U_j^0)$ 이다($s_{\phi,j}$=방전 방향 부호 인자, \S\ref{sec:notation};
식~\eqref{eq:Aj} 와 같은 sign convention). ideal 극한($\Omega_j=0$)에서 $\mu^0+RT\ln[\xi_\eq/(1-\xi_\eq)]
=s_{\phi,j}F(V_n-U_j^0)$ 이고, 상수 $\mu^0$ 를 기준 전위에 흡수해 $U_j\equiv U_j^0-\mu^0/(s_{\phi,j}F)$ 로
재정의하면
\begin{equation}
RT\ln\frac{\xi_\eq}{1-\xi_\eq}=s_{\phi,j}F(V_n-U_j).
\label{eq:eqbal_ideal}
\end{equation}

\textbf{(iv) $\xi_\eq$ 에 대해 풀기(한 줄씩).} 식~\eqref{eq:eqbal_ideal} 양변을 $RT$ 로 나누고 폭
$w_j\equiv RT/F$ 를 정의하면 $\ln[\xi_\eq/(1-\xi_\eq)]=s_{\phi,j}(V_n-U_j)/w_j$. 양변에 지수를 취해
$\dfrac{\xi_\eq}{1-\xi_\eq}=E$, $E\equiv\exp[s_{\phi,j}(V_n-U_j)/w_j]$. 이를 풀면 $\xi_\eq=E(1-\xi_\eq)$
$\Rightarrow\xi_\eq(1+E)=E\Rightarrow\xi_\eq=E/(1+E)$, 분자·분모를 $E$ 로 나눠 $1/E=\exp[-s_{\phi,j}(V_n-U_j)/w_j]$
이므로
\begin{equation}
\boxed{\;\xi_{\eq,j}(V_n,T)=\frac{1}{1+\exp[-s_{\phi,j}(V_n-U_j(T))/w_j(T)]}\;,\qquad w_j=\frac{RT}{F}\ (\text{ideal})\;.}
\label{eq:logistic}
\end{equation}
모든 중간 단계가 위에 명시됐다(방전 기준 $s_{\phi,j}=+1$ 이 기본이며, 이때 $V_n$ 증가가 $\xi_\eq$ 를 증가시킨다).
$\Omega_j\ne0$ 이면 식~\eqref{eq:eqbal_ideal} 우변에 $-\Omega_j(1-2\xi_\eq)$ 가 더해져 \emph{음함수}(implicit
isotherm)가 되며 단순 logistic 이 아니다 — 이때 $w_j$ 를 effective width 로 두는 것은 닫힌 유도가 아니라
coarse-grained ansatz 임을 명시한다(\AL{2} BOUNDED; $\Omega_j$ 가 크면 다중해·상분리 가능). 이는 매끄러운 logistic 곡선이며 $0\!\to\!1$ 급점프가
아니다. $\Omega_j\ne0$ 이면 폭·형태가 매끄럽게 보정되며, 실제로는 nonideality$\cdot$도메인 불균일$\cdot$유한
전이폭을 흡수하는 effective width $w_j$ 로 두어 $RT/F$ 에 강제로 묶지 않는다(\AL{4}). 이때
$n_j^{\eff}\equiv RT/(Fw_j)$ 가 effective 전자수 역할을 하며 Ch.3$\cdot$Ch.4 로 전달된다.

\textbf{저온 거동(서 단계 (1) 의 근거; 무비약).} logistic 의 peak 형상은 그 도함수로 결정된다. 식~
\eqref{eq:logistic} 에서 $z\equiv s_{\phi,j}(V_n-U_j)/w_j$ 로 두면 $\xi_\eq=(1+e^{-z})^{-1}$ 이고
$\dd z/\dd V_n=s_{\phi,j}/w_j$. 연쇄법칙으로 $\partial\xi_\eq/\partial V_n=(\dd\xi_\eq/\dd z)(\dd z/\dd V_n)$
이며, $\dd\xi_\eq/\dd z=-(1+e^{-z})^{-2}\cdot(-e^{-z})=e^{-z}/(1+e^{-z})^2$ 다. 여기서
$1-\xi_\eq=e^{-z}/(1+e^{-z})$ 이므로 $e^{-z}/(1+e^{-z})^2=[1/(1+e^{-z})]\cdot[e^{-z}/(1+e^{-z})]
=\xi_\eq(1-\xi_\eq)$ (표준 logistic 항등식)이다. 따라서
\begin{equation}
\frac{\partial\xi_{\eq,j}}{\partial V_n}=\frac{s_{\phi,j}\,\xi_{\eq,j}(1-\xi_{\eq,j})}{w_j(T)}
\label{eq:dxidV}
\end{equation}
로 결정되어 peak 높이 $\propto 1/w_j$, 폭 $\propto w_j$ 다($s_{\phi,j}=+1$ 방전 기준). 저온이면 \emph{ideal}
$w_j=RT/F$ 가 \emph{감소}하므로
peak 가 더 좁고 가팔라진다 — 즉 꼬리가 더 빨리 끝난다. 수치로 $25^\circ$C 에서 $w=RT/F\approx25.7$ mV,
$-15^\circ$C 에서 $\approx22.2$ mV($8.314\times298/96485=25.68$ mV, $\times258/96485=22.23$ mV).
\textbf{따라서 평형 열역학의 예측은 ``저온 $\to$ 짧은 꼬리''} 이고, 관찰(저온 긴 꼬리)은 평형 broadening 으로
설명할 수 없다. 꼬리의 주역은 평형이 아니라 동역학이어야 한다(\S\ref{sec:rate}--\ref{sec:dyn}).
\begin{flagbox}
\textbf{평형 형태는 데이터가 정한다.} logistic 대신 erf/Gaussian-cumulative isotherm 을 쓰는 것은 흑연에
대한 열역학적 유도가 문헌에 없고(무질서 반도체 Gaussian-disorder 모형 \cite{baessler1993} 과의 유추 +
경험적 peak-fit 수준) FLAGGED ansatz 다. 실제 형태는 저속 OCV 의 near-peak 감쇠를 $\log(dQ/dV)$ 대
$(V-U)$(지수=logistic) 또는 $(V-U)^2$(가우시안)로 보고 판별한다(\S\ref{sec:falsify}). \textbf{본 장의
꼬리 결론은 평형 형태에 무관}하다 — post-peak 에서 $\dd\xi_\eq/\dd q\simeq0$ 만 쓰기 때문이다.
\end{flagbox}

% ====================================================================
\section{주역 (I): 극판 전위에 의한 배리어 낮춤과 속도상수}\label{sec:rate}
% ====================================================================
이 절이 본 장의 핵심이다 — \emph{극판 전위가 어떻게 상변이 속도를 바꾸는가}.

\textbf{세 전위의 구분.} 측정되는 apparent 전위는 내부 전위에 분극이 더해진 값이다(\AL{6}):
\begin{equation}
V_{n,\app}=V_n+s_I|I|R_n(q,T,|I|).
\label{eq:Vapp}
\end{equation}
상변이 속도를 정하는 구동 전위는 $V_{n,\drive}$ 이며 가장 보수적인 기본형은 $V_{n,\drive}=V_n$ 이다(국소
계면 과전압을 분리하는 일반형 $V_{n,\drive}=V_n+\eta_{n,\drive}$ 는 Ch.3). \textbf{$R_n$(전압축 이동)과
아래의 속도 보조항을 같은 데이터로 동시에 자유 피팅하지 않는다} — 역할이 다르다($R_n$=전압 강하,
$\chi_j$=상변이 가속). 이 구분이 무너지면 둘이 서로를 흉내 내어 식별이 불가능해진다(\S\ref{sec:falsify}).

\textbf{구동력.} 한 전이의 구동력은 구동 전위가 전이 중심 전위에서 벗어난 정도다(\AL{11}):
\begin{equation}
\mathcal A_j=s_{\phi,j}F\,[V_{n,\drive}-U_j(T)].
\label{eq:Aj}
\end{equation}
$F[\mathrm V]$ 가 $\mathrm{J/mol}$ 차원이므로 $\mathcal A_j$ 는 몰당 자유에너지 구동력이다. 부호 인자
$s_{\phi,j}$ 는 선택한 방전 방향 전이의 구동력 부호를 정한다: $\chi_j\ge0$ 이고 $s_{\phi,j}[V_{n,\drive}-U_j]>0$
일 때 방전 방향 전이가 촉진된다. (본 장은 ideal 폭 $n_j^{\eff}=RT/(Fw_j)=1$ 을 기본으로 두므로 위 형태를 쓴다;
일반형 $\mathcal A_j=s_{\phi,j}\,n_j^{\eff}F(V_{n,\drive}-U_j)$ 는 effective width 와 함께 Ch.3 에서 쓴다.)

\textbf{유효 배리어(무비약).} 고유 활성화 자유에너지는 전이상태 이론(\AL{10}, \cite{eyring1935})으로
$\Delta G_{a,j}(T)=\Delta H_{a,j}-T\Delta S_{a,j}$ 다. 전위 구동력이 이 배리어를 \emph{선형으로} 낮춘다 —
이것이 Brønsted--Evans--Polanyi/Butler--Volmer 류의 표준 자유에너지 관계다(\AL{11},
\cite{evanspolanyi1938,bardfaulkner,bazant2013}).

\emph{선형의 출처(소구동력 1차 Taylor; 무생략).} Marcus 이론에서 구동력 $\mathcal A_j$ 가 주어졌을 때
활성화 장벽은 포물면 교차로부터 $\Delta G^\ddagger(\mathcal A_j)=\dfrac{(\lambda-\mathcal A_j)^2}{4\lambda}
=\dfrac{\lambda}{4}-\dfrac{\mathcal A_j}{2}+\dfrac{\mathcal A_j^2}{4\lambda}$ ($\lambda$=재구성 에너지)다.
$\mathcal A_j=0$ 둘레 1차 Taylor 전개를 하면 상수항 $\Delta G^\ddagger(0)=\lambda/4$ 와 1차 항
$\big[\dd\Delta G^\ddagger/\dd\mathcal A_j\big]_{0}\mathcal A_j=-\tfrac12\mathcal A_j$ 가 남고, 2차 항
$\mathcal A_j^2/(4\lambda)$ 은 $|\mathcal A_j|\ll\lambda$ 에서 무시된다. 따라서 1차 계수의 크기를 본 장의 linear sensitivity
$\chi_j$ 로 둔다. \emph{대칭 Marcus 포물면(forward/backward 동일 곡률 $\lambda$)의 1차 전개는 정확히 $\chi_j=1/2$
단일값만 준다}; $\chi_j\ne1/2$(일반 $0\le\chi_j\le1$)은 1차 Taylor 자체에서 나오는 것이 아니라 \emph{곡률 비대칭}
($\lambda_f\ne\lambda_b$) 또는 BEP/Butler--Volmer 선형 자유에너지 관계의 경험적 기울기에서 오며
(\AL{11}, \cite{evanspolanyi1938,bardfaulkner}), 그 일반 범위로 둔다. 이 1차 항이 장벽을 낮추므로($\mathcal A_j>0$ 일 때 $-\chi_j\mathcal A_j<0$)
\begin{equation}
\boxed{\;\Delta G_{\eff,j}=\Delta G_{a,j}(T)-\chi_j\,\mathcal A_j\;=\;G-\chi_j\mathcal A_j\;}
\label{eq:Geff}
\end{equation}
($G\equiv\Delta G_{a,j}$ 를 mode 의 intrinsic barrier 로 표기, 위 $\lambda/4$ 등 상수항 흡수;
\S\ref{sec:spectrum} 에서 분포). 곡률(2차 항)이 중요해지는 $|\mathcal A_j|\sim\lambda$ 영역은 아래 Marcus
bound 로 분리한다. \emph{주의}: 이 $\chi_j\mathcal A_j$ 는 forward/backward 를 같은 배율로 키우는
\emph{방향 없는 common-mode mobility 가속}이며(keystone 참조), forward 장벽만 낮추는 표준 Butler--Volmer
와는 다르다 — 방향성 분해는 Ch.3 Level B 에서 다룬다. 구동력이
클수록($\mathcal A_j>0$) 유효 배리어가 낮아진다. 속도상수는 Eyring rate(\AL{10})
\begin{equation}
\boxed{\;k_j=k_0(T)\exp\!\Big[-\frac{\Delta G_{\eff,j}}{RT}\Big]\;,\qquad k_0(T)=\frac{k_BT}{h}\kappa(T)\;.}
\label{eq:kact}
\end{equation}
학부 수준에서는 Boltzmann 인자 $k_j\propto e^{-\Delta G_{\eff,j}/RT}$ 만으로 이후 논증에 충분하다 — 본 장의
모든 결론이 속도의 \emph{비}로 정의되는 길이 $L=|I|/(Q_\cell k_j)$(\S\ref{sec:dyn})로만 쓰여 prefactor 절대값
$k_0$ 가 상쇄되기 때문이다. prefactor 가 $k_BT/h$ 로 \emph{고정}(임의 상수 아님)이라는 절대값 자체는 전이상태
이론(\cite{eyring1935}) 결과로 \emph{인용}한다. $\Delta G_{\eff,j}\le0$
(강한 전위 구동)은 선형 근사가 강구동 영역에 들어간 것으로 표시하되 $\max(\cdot,0)$ 같은 절단으로 자르지
\emph{않는다}(GS-4; 절단이 필요해 보이면 그것은 아래 물리적 병목으로 다룬다).
\begin{keybox}
\textbf{$\chi_j\mathcal A_j$ 는 ``방향 있는 배리어 낮춤''이 아니라 ``방향 없는 완화 속도(mobility) 가속''이다
(Level A).} 무비약 근거: 2-상태 반응속도론에서 net rate $\dot\xi_j=r_j^+(1-\xi_j)-r_j^-\xi_j$. 정상점은
$\dot\xi_j=0$ 에서 $r_j^+(1-\xi_{\eq,j})=r_j^-\xi_{\eq,j}\Rightarrow r_j^+=(r_j^++r_j^-)\xi_{\eq,j}
\Rightarrow\xi_{\eq,j}=r_j^+/(r_j^++r_j^-)$ 다. $r_j^\pm$ 를 $\xi_j$ 의 작은 변화에 대해 국소 상수로 두면
net rate 를 직접 재배열할 수 있다: $\dot\xi_j=r_j^+-r_j^+\xi_j-r_j^-\xi_j=r_j^+-(r_j^++r_j^-)\xi_j
=(r_j^++r_j^-)\big[r_j^+/(r_j^++r_j^-)-\xi_j\big]=(r_j^++r_j^-)(\xi_{\eq,j}-\xi_j)$. 따라서
$\dot\xi_j=k_j(\xi_{\eq,j}-\xi_j)$, $k_j=r_j^++r_j^-$ 가 인수분해로 동시에 나온다(\AL{12}). \emph{강조}: 비약은
``$r_j^\pm$ 를 ($\xi_j$ 에 무관한) 국소 상수로 본다''는 \emph{가정}에만 있고, 그 가정 위에서는 위 재배열이 Taylor
전개가 아니라 대수적 항등이다(반대로 $r_j^\pm(\xi_j)$ 면 같은 식이 국소 선형화가 된다). 즉 $k_j$=완화 속도(목표에 접근하는 빠르기),
$\xi_{\eq,j}$=목표(방향). 식~\eqref{eq:kact} 처럼 $\chi_j\mathcal A_j$ 가 \emph{$k_j$
전체}에 들어가면 $r_j^+,r_j^-$ 가 같은 배율로 커져 비 $r_j^+/r_j^-=\xi_\eq/(1-\xi_\eq)$ 가 \emph{불변}
— 즉 방향(목표)은 열역학(\S\ref{sec:eq})이 전담하고, $\chi_j\mathcal A_j$ 는 방향 없는 속도 scale 일 뿐이다.
\emph{방향을 가진} 배리어 낮춤(forward 만 $-\beta\mathcal A$, backward 는 $+(1-\beta)\mathcal A$ 라 비가 바뀜)은 \textbf{Ch.3 의 Level B} 에서 도입한다.
그때는 별도의 transfer coefficient $\beta_j$ 를 정의하고, 본 장의 $\chi_j$ 와의 대응은 추가 kinetic model 아래에서만 검토한다(\AL{11}).
\textbf{한정}: 본 장 Level A 의 ``공통 배율''
(forward/backward 동일 가속)은 Marcus 의 \emph{비대칭} 장벽 이동($\lambda\mp\mathcal A_j$)을 대칭으로 평균낸
\emph{모형 선택}이며, 이 대칭 극한에서만 forward/backward 가속이 세부 균형(detailed balance) 비를 보존한다. Level B(Ch.3)와
일치하는 것은 \emph{정상 target 극한} $\xi_{\mathrm{ss},j}\to\xi_{\eq,j}$ 한정이다(이것이 Ch.3 연결 조건이다). 본 장
결론식에서는 ``배리어 낮춤''이라는 방향성 함의 용어를 피한다.
\end{keybox}
\begin{boundbox}
\textbf{Marcus bound(\AL{15}).} 선형 $-\chi_j\mathcal A_j$ 는 작은 구동력의 1차 근사다. Marcus 이론
\cite{marcus1956} 은 곡률을 주므로 $|\mathcal A_j|$ 가 클 때(역전 영역 포함) 깨진다. 본 선형식은 꼬리 영역
$V_{n,\drive}\approx U_j$ 에서 유효하며 전역 정당화가 아니다.
\end{boundbox}

\textbf{물리적 병목(속도 절단이 아니라 직렬 합성).} 전하전달이 빨라져도 다음 단계(전해질/고체 수송, 유한
활성 자리, 고체상 확산, 유한 전류)가 전체 진행을 제한할 수 있다. 이를 \emph{시간척도의 직렬 합}으로 둔다
(수치 cap 이 아니라 율속 단계 합성):
\begin{equation}
\frac{1}{k_j^{\eff}}=\frac{1}{k_j}+\frac{1}{k_{j,\lim}},
\label{eq:kphys}
\end{equation}
$k_{j,\lim}\to\infty$ 면 활성화 지배($k_j^{\eff}\to k_j$), 반대면 병목 지배로 매끄럽게 환원된다.\footnote{%
$k_{j,\lim}$ 은 수송$\cdot$유한 자리$\cdot$고체 확산$\cdot$유한 전류 등 병목의 합($1/k_{j,\lim}=\sum 1/k_{j,m}$)
이다. 독립 자료 없이 자유 파라미터로 추가하면 $k_0,\Delta G_a,\chi_j$ 와 강하게 상관되므로 물리적 근거와
독립 제약이 있을 때만 쓴다. 유한 전류를 매끄럽지 않은 clamp($\min$/step)로 강제하지 않는다(GS-4); 그런
편의 처리와 피팅 안정화는 부록 \S\ref{sec:ch1_numeric} 에서 다룬다.} 이하 $k_j$ 는 이 $k_j^{\eff}$ 를 가리킨다.

% ====================================================================
\section{주역 (II): 진행률 동역학과 single-mode kernel}\label{sec:dyn}
% ====================================================================
상변이는 평형 진행률을 향해 유한 속도로 완화된다(\AL{12}, \cite{onsager1931,degrootmazur1962}):
\begin{equation}
\boxed{\;\frac{\dd\xi_j}{\dd t}=k_j\,[\xi_{\eq,j}(V_n,T)-\xi_j]\;.}
\label{eq:dxidt}
\end{equation}
$\xi_j$ 가 평형을 즉시 따라가지 못하면 전압 곡선과 ICA/DVA 에 꼬리가 생긴다. 정전류 구간($|I|>0$)에서는
$\dd q/\dd t=|I|/Q_\cell$ 이므로 그 역수가 $\dd t/\dd q=Q_\cell/|I|$ 이고, 연쇄법칙
$\dd\xi_j/\dd q=(\dd\xi_j/\dd t)(\dd t/\dd q)$ 에 식~\eqref{eq:dxidt} 의 $\dd\xi_j/\dd t=k_j[\xi_{\eq,j}-\xi_j]$
를 넣으면
\begin{equation}
\frac{\dd\xi_j}{\dd q}=\frac{Q_\cell}{|I|}\,k_j\,[\xi_{\eq,j}-\xi_j].
\label{eq:dxidq}
\end{equation}

\textbf{Single-mode kernel.} 한 mode 의 지연을 $r=\xi_{\eq,j}-\xi_j$ 라 하자. peak 를 지난 영역(post-peak)
에서 목표가 거의 정지하고($\dd\xi_{\eq,j}/\dd q\simeq0$) 구동력도 완만하면($\dd\mathcal A_j/\dd q\simeq0$,
즉 $V_n$ 완만 — \emph{별개 가정 두 개}임에 유의), 식~\eqref{eq:dxidq} 는 1계 선형 ODE 가 된다.
\emph{무비약 유도}: 정의 $r=\xi_{\eq,j}-\xi_j$ 를 $q$ 로 미분하면 $\dd r/\dd q=\dd\xi_{\eq,j}/\dd q-\dd\xi_j/\dd q$
이고, post-peak 가정 $\dd\xi_{\eq,j}/\dd q\simeq0$ 으로 첫 항이 사라져 $\dd r/\dd q=-\dd\xi_j/\dd q$ 다. 여기에
식~\eqref{eq:dxidq} $\dd\xi_j/\dd q=(Q_\cell k_j/|I|)(\xi_{\eq,j}-\xi_j)=(Q_\cell k_j/|I|)\,r$ 를 넣고
$L\equiv|I|/(Q_\cell k_j)$ 로 두면 $\dd r/\dd q=-(1/L)\,r$ 다. 여기서 $L$ 은 일반적으로 $q$ 의 함수다($k_j$ 가
구동력 $\mathcal A_j$ 에 지수 의존, 식~\eqref{eq:kact}). \emph{그러나} 위 가정 (나) $\dd\mathcal A_j/\dd q\simeq0$
이면 $k_j$, 따라서 $L$ 이 이 구간에서 $q$-무관(상수)이므로 $L$ 을 적분 밖으로 뺄 수 있다 — 이것이 ``$L$ 일정''의
근거다(별개 가정이 아니라 (나)의 \emph{귀결}). 변수분리 $\dd r/r=-\dd q/L$ 를 $q_a$ 에서 $q$ 까지 적분하면
$\int_{q_a}^q \dd q'/L=(q-q_a)/L$, 즉 $\ln[r(q)/r(q_a)]=-(q-q_a)/L$ 이고 양변에 지수를 취해 지수 감쇠가 따라온다:
\begin{equation}
r(q)\simeq r(q_a)\exp\!\Big[-\frac{q-q_a}{L}\Big],\qquad \boxed{\;L=\frac{|I|}{Q_\cell\,k_j}\;}
\label{eq:single_kernel}
\end{equation}
(차원: $\mathrm A=\mathrm{C/s}$, 속도상수 $k$ 는 $\mathrm{s^{-1}}$, $Q_\cell$ 은 $\mathrm C$ 이므로
$|I|/(Q_\cell k)=(\mathrm{C/s})/(\mathrm C\cdot\mathrm{s^{-1}})=(\mathrm{C/s})/(\mathrm{C/s})=1$, 즉 무차원).
이것은 \emph{관측 꼬리 전체가 단일
지수}라는 뜻이 아니라 \textbf{하나의 mode 가 만드는 지수 kernel} 이다. 관측 꼬리는 여러 mode 의 중첩이다
(\S\ref{sec:kernel}).

\textbf{잔차에서 진행률 도함수로(무비약 연결).} 다음 절의 중첩은 잔차 $r$ 이 아니라 \emph{진행률 도함수}
$\dd\xi_j/\dd q$ 의 합이다. 둘의 연결은 식~\eqref{eq:dxidq} 에서 $\dd\xi_j/\dd q=(Q_\cell k_j/|I|)\,r=r/L$
이므로
\begin{equation}
\frac{\dd\xi_j}{\dd q}=\frac{r(q_a)}{L}\exp\!\Big[-\frac{q-q_a}{L}\Big].
\label{eq:single_dxidq}
\end{equation}
바로 이 \textbf{$1/L$ 인자}가 \S\ref{sec:kernel} kernel 의 $1/L$ 출처다 — 중첩은 각 mode 의 \emph{진행률
기여}를 합한 것이지 잔차를 합한 것이 아니다.

\textbf{저온 거동(서 단계 (2)).} $k_j=k_0e^{-(G-\chi\mathcal A)/RT}$ 가 저온서 작아져 $L$ 이 커진다 $\Rightarrow$
꼬리가 길어진다 — 관찰과 일치. 빠른 kinetics 극한($k_j\to\infty$)에서는 $L\to0$ 이며 $\dd\xi_j/\dd q\to
\dd\xi_{\eq,j}/\dd q$ 로 접근한다(작은 지연 $\times$ 큰 속도의 곱이 유한; bracket 이 단순히 0 이 되는 것이
아니라 점근 균형의 결과다).

% ====================================================================
\section{주역 (III): 배리어 분포 $\to$ relaxation-length spectrum}\label{sec:spectrum}
% ====================================================================
\textbf{물리.} 실제 전극은 단일 배리어 $G$ 가 아니라 입자 크기$\cdot$국소 staging 환경$\cdot$결정립계$\cdot$결함
$\cdot$응력 차이로 인한 \emph{배리어 분포} $\rho_G(G;T)$ 를 가진다(\AL{13}). 진행률을 특정 배리어 기준으로
자르지 않고, 배리어 $g$ 를 가진 집단의 진행률을 $\xi_j(g,t)$ 로 두어 분포로 평균한다($\int_0^\infty
\rho_G(g;T)\,\dd g=1$):
\begin{equation}
\xi_j(t)=\int_0^\infty \rho_G(g;T)\,\xi_j(g,t)\,\dd g,\qquad
\frac{\dd\xi_j(g,t)}{\dd t}=k_j(g)\,[\xi_{\eq,j}-\xi_j(g,t)].
\label{eq:xi_dist}
\end{equation}
\begin{flagbox}
\textbf{독립 완화 가정(평균장; \AL{13}).} 식~\eqref{eq:xi_dist} 는 배리어 $g$ 별 집단이 \emph{서로 독립적으로}
1차 완화한다고 본다(mode 간 결합 무시). 그러나 흑연 staging 은 nucleation$\cdot$상경계 이동을 동반하는
\emph{협동적 1차 상전이}다(\AL{3}, \cite{funabiki1999ea,funabiki1999jes}). 따라서 ``독립 평행 1차 완화의 분포''
근사는 mode 간 상호작용(상경계 결합$\cdot$Avrami 형 비선형 동역학)을 무시한 \emph{평균장 ansatz} 이며, 그 타당성은
FLAGGED — Ch.3(반응속도론 일반화)$\cdot$본 장 부록 \S\ref{sec:ch1_numeric}(수치)에서 검증한다.
\end{flagbox}

\textbf{배리어 $\to$ 길이 지수 매핑(꼬리가 늘어지는 근원).} 관측 꼬리의 모양은 $\rho_G$ 와 \emph{같지 않다}
— 배리어 $G$ 가 먼저 속도로, 다시 길이로 \emph{지수 변환}되기 때문이다. 식~\eqref{eq:single_kernel} 의
$L=|I|/(Q_\cell k_j)$ 에 식~\eqref{eq:kact} $k_j=k_0(T)\exp[-(G-\chi_j\mathcal A_j)/RT]$(mode 별 intrinsic
barrier $G\equiv\Delta G_{a,j}$, $\Delta G_{\eff,j}=G-\chi_j\mathcal A_j$)를 대입한다. \emph{이 닫힌 지수 매핑은
활성화 지배 극한}($k_j^{\eff}\simeq k_j$, 즉 식~\eqref{eq:kphys} 의 $k_{j,\lim}\to\infty$)\emph{을 전제}한다 — 꼬리
영역($V_{n,\drive}\approx U_j$, 저구동)에서는 활성화가 율속이라 정당하나, 병목이 유효하면
$1/k_j^{\eff}=1/k_j+1/k_{j,\lim}$ 가 지수 밖 유리식으로 들어가 $L(G)$ 가 단순 지수형이 아니다(그 경우의 수치 풀이는
부록 \S\ref{sec:ch1_numeric}). $L$ 은 $k_j$ 의 \emph{역수}에 비례하므로 $1/k_j=(1/k_0)\exp[+(G-\chi_j\mathcal A_j)/RT]$ 로
\emph{지수 부호가 반전}되어(속도가 작을수록 길이가 길다)
\begin{equation}
\boxed{\;L(G,T,V_{n,\drive})=\frac{|I|}{Q_\cell\,k_0(T)}\exp\!\Big[\frac{G-\chi_j\mathcal A_j}{RT}\Big]\;.}
\label{eq:L_of_G}
\end{equation}
양변에 로그를 취하면 $\ln L=\ln L_0+(G-\chi_j\mathcal A_j)/RT$ 이다. 즉 배리어의 표준편차 $\sigma_G$($\rho_G$ 의
표준편차)가 $\ln L$ 에서 $\sigma_G/RT$ 로 나타나, \textbf{저온일수록 $\ln L$ 의 퍼짐이 커진다} — 작은 배리어 분산이 길이에서
지수적으로 확대된다. 이것이 ``저온 긴 꼬리''의 근원이며, 고정된 배리어 분포가 온도를 통해 stretched 거동을
만든다는 메커니즘은 무질서계$\cdot$빠른 이온 전도체 동역학에서 확립돼 있다(\AL{14},
\cite{svare2000,johnston2006,lindsey1980}; 단 graphite ICA 꼬리로의 \emph{적용}은 BOUNDED, \S\ref{sec:kernel}).

\textbf{변수변환: 배리어 분포 $\to$ 완화 길이 분포(무비약, 단계별).} 우리가 만들 변수 $A_L(L)$ 은 ``완화 길이가
$L$ 인 mode 가 꼬리에 \emph{얼마나 기여}하는가''의 분포다. 이는 \emph{이미 아는} 배리어 분포 $\rho_G(G)$ 를
식~\eqref{eq:L_of_G} 의 $L(G)$ 로 \emph{좌표만 바꾼} 것 — 새 물리를 넣는 게 아니라 같은 mode 집단을 길이 축에서
다시 세는 것이다. 세 단계로 만든다.

\emph{단계 1 — 자코비안(밀도 보존).} 확률(mode 분율)은 좌표를 바꿔도 보존되므로 $A_L^{\mathrm{prob}}(L)\,\dd L
=\rho_G(G)\,\dd G$. 역변환(식~\eqref{eq:L_of_G} 을 $G$ 에 대해 풀면) $G(L)=\chi_j\mathcal A_j+RT\ln(L/L_0)$
($L_0\equiv|I|/(Q_\cell k_0)$; $G=0$ 일 때 $L=L_{\min}$, 단계 2 에서 정식화)에서 $\dd G/\dd L=RT/L>0$ 이라
일대일이므로
\begin{equation}
A_L^{\mathrm{prob}}(L)=\rho_G\big(G(L)\big)\,\Big|\frac{\dd G}{\dd L}\Big|=\rho_G\big(G(L)\big)\,\frac{RT}{L}.
\label{eq:spectrum_prob}
\end{equation}
\emph{의미}: $A_L^{\mathrm{prob}}(L)\dd L$ $=$ 완화 길이가 $[L,L+\dd L]$ 인 mode 의 \emph{분율}, 규격 보존
($\int A_L^{\mathrm{prob}}\dd L=\int\rho_G\dd G=1$). 자코비안 $RT/L$ 은 ``배리어 한 칸 $=$ 길이 몇 칸''의 환산율일
뿐 임의 가중이 아니다.

\emph{단계 2 — 지지 범위(음의 배리어 금지).} 활성화 배리어는 음수가 못 되므로($G\ge0$, \S\ref{sec:notation})
식~\eqref{eq:L_of_G} 의 단조성상 $G\ge0\Leftrightarrow L\ge L_{\min}$, $L_{\min}\equiv L_0 e^{-\chi_j\mathcal A_j/RT}$
($G=0$ 대응 최단 길이)다. 따라서 length-domain은 $L\in[L_{\min},\infty)$로 제한된다. 이것은 artificial switching law가 아니라
음의 activation barrier를 적분하지 않기 위한 domain condition이다.

\emph{단계 3 — 개수 밀도 $\to$ 용량 가중(측도 변환, 무비약).} 단계 1$\cdot$2 의 $A_L^{\mathrm{prob}}$ 는 ``각 길이에
mode 가 몇이나''(개수 밀도)다. 그러나 관측 꼬리는 \emph{용량 가중} 진행 $\Theta=Q_p^{-1}\sum_m Q_m\xi_m$
(\S\ref{sec:fiteq}; $Q_m$=mode 의 용량 기여)로 나타나므로, 길이별 기여는 \emph{개수}$\times$\emph{mode 당 용량}이다.
이산 합을 연속 적분으로 옮기면
\begin{equation*}
\Theta=Q_p^{-1}\sum_m Q_m\,\xi_m\;\xrightarrow{\text{연속극한}}\;Q_p^{-1}\!\int Q(L)\,\xi(L)\,A_L^{\mathrm{prob}}(L)\,\dd L,
\end{equation*}
즉 개수 밀도 $A_L^{\mathrm{prob}}$ 에 \emph{mode 당 용량} $Q(L)$ 이 곱해진다. 이 용량 가중을 (무차원화한)
$A_0(L)\equiv Q(L)/\bar Q$ 로 두면 $A_0$ 는 개수 밀도 $\rho_G$ 에서 \emph{용량 측도}로 가는 Radon--Nikodym
가중이다 — 입자 크기 분포$\cdot$활성 site 밀도 같은 독립 물리량에서 오며 임의 피팅 knob 이 아니다. 관측 spectrum 은
그 곱이다:
\begin{equation}
\boxed{\;A_L(L)=A_0(L)\;\underbrace{\rho_G\big(G(L)\big)\,\frac{RT}{L}}_{\displaystyle A_L^{\mathrm{prob}}(L)\ (\text{규격 보존})},\qquad L\ge L_{\min}.\;}
\label{eq:spectrum}
\end{equation}
$A_0$ 는 \emph{임의 피팅 knob 이 아니라} 독립 물리량(입자 크기 분포 등)에서 와야 한다. $A_0\equiv1$ 이면
$A_L=A_L^{\mathrm{prob}}$ 로 규격 보존, $A_0\ne1$ 이면 $A_L$ 은 규격화 안 된 \emph{진폭 스펙트럼}이다 — \textbf{이하
kernel 은 이 $A_L$(진폭 spectrum)을 쓴다}(관측 꼬리는 진폭 가중 중첩이므로). 자코비안 $RT/L$ 과 아래 kernel 의
$1/L$ 인자는 \emph{출처가 다른 별개}임에 유의(\S\ref{sec:kernel}).
\begin{boundbox}
\textbf{비유일성(\AL{16}).} stretched 거동에서 $\rho_G$ 로의 역매핑은 \emph{비유일}하다(재포획 등 다른
메커니즘도 같은 꼬리를 낼 수 있다 \cite{baessler1993}). 따라서 $\rho_G$ 를 관측 꼬리에서 \emph{역산하지
않고}, 독립 측정/모형한 $\rho_G$ 로부터 꼬리를 \emph{forward 예측}만 한다(\S\ref{sec:falsify}).
\end{boundbox}

% ====================================================================
\section{관측 꼬리 = kernel integral}\label{sec:kernel}
% ====================================================================
\textbf{$A_L$ 의 물리 의미.} $A_L(L)\,\dd L$ 은 \emph{완화 길이가 $[L,L+\dd L]$ 인 mode 들이 전체 진행에 기여하는
분율}이다 — 즉 관측 꼬리를 ``완화 길이별로 분해한 분포''다. 빠른 mode(작은 $L$)는 peak 근처에서 빨리 끝나고,
느린 mode(큰 $L$)가 꼬리를 멀리 끌고 간다. 따라서 관측 꼬리는 \emph{서로 다른 $L$ 의 지수 감쇠를 $A_L$ 로 가중
합한 것}이다: 전체 진행률의 post-peak 꼬리는 각 mode 의 진행률 기여 $\dd\xi_j/\dd q=(r(q_a;L)/L)e^{-(q-q_a)/L}$
(식~\eqref{eq:single_dxidq})를 $A_L$ 로 가중 합한 것이다. 여기서 mode 별 초기 잔차 $r(q_a;L)$ 를 post-peak
공통값 $\bar r$ 로 \emph{근사}하면 $\bar r$ 가 적분 밖으로 빠지고, 이를 진폭 규격에 재흡수($A_L\!\leftarrow\!\bar r\,A_L$)해
아래 식을 얻는다 — $r(q_a)$ 가 사라진 게 아니라 $A_L$ 의 크기 인자로 흡수된 것이다:
\begin{equation}
\boxed{\;\frac{\dd\Theta_\mathrm{tail}}{\dd q}\simeq\int_{L_{\min}}^\infty A_L(L)\,\frac1L\,
\exp\!\Big[-\frac{q-q_a}{L}\Big]\,\dd L\;.}
\label{eq:kernel_integral}
\end{equation}
이 균일 근사는 BOUNDED 다: 빠른 mode(작은 $L$)는 $q_a$ 에서 이미 더 진행해 $r(q_a;L)$ 가 작을 수 있어 원칙적으로
$L$-의존이다. $r(q_a;L)$ 이 $L$ 의 \emph{감소}함수이면 작은-$L$ 기여를 과대평가해 \emph{꼬리 앞부분이 실제보다
가팔라지는} 쪽으로 편향되며, 이 편향은 $r(q_a;L)$ 를 직접 보유하는 Plan B(g-grid, \S\ref{sec:planB})와 대조해
정량화한다. 피적분의 $1/L$ 인자는 진행률 도함수(식~\eqref{eq:single_dxidq})에서 온 것이고, spectrum 변환의
자코비안 $RT/L$ 은 $A_L$ 안에 이미 들어 있어 둘은 \emph{별개} 출처다 — 임의 가중이 아니다. $A_L$ 은 \S\ref{sec:spectrum}
에서 $L\ge L_{\min}>0$로 정의되므로 적분의 실효 하한은 $L_{\min}$이고 $1/L$ 의 $L\!\to\!0$ 발산은 일어나지 않는다.
이 적분은 relaxation-rate($1/L$) 분포 위의 Laplace 형 변환이며, 단일 지수의 \emph{연속 중첩}이 일반적으로 비지수
(stretched)$\sim$power-law 꼬리를 준다(\AL{14}, \cite{johnston2006}; KWW stretched 함수와 완화시간 분포의
등가성 \cite{lindsey1980}). single-mode(식~\eqref{eq:single_kernel})는 amplitude-bearing form $A_L(L)=\Theta_0\delta(L-L^\ast)$ 인 특수해다
($\Theta_0$=post-peak residual amplitude, $L^\ast$=그 mode 의 실제 완화 길이 $|I|/(Q_\cell k_j)$; \S\ref{sec:spectrum} 의 prefactor 길이 $L_0$ 와 구별).
목표가 움직이는 일반(자기참조 포함) 경우는 동일 kernel 을 \emph{목표} $\Theta_\eq$ 에 작용시킨 \S\ref{sec:volterra}
(식~\eqref{eq:volterra})로 일반화된다(잔차가 아니라 목표에 곱함 — \S\ref{sec:volterra} 의 이중계상 경고 참조).
\begin{boundbox}
\textbf{유효범위(\AL{14}, 정직 표기).} ``고정 배리어 분포 $\to$ 지수 매핑 $\to$ 저온 긴 꼬리''의 메커니즘은
빠른 이온 전도체$\cdot$무질서 완화에서 확립돼 있다(\cite{svare2000,johnston2006,lindsey1980}). 다만 (i)
특정 $\rho_G$ 형태가 \emph{정확히 어떤} stretched 지수를 주는지는 데이터/수치(부록 \S\ref{sec:ch1_numeric})로 정하며, (ii) 위
문헌은 전도도/유전 완화 맥락이라 graphite ICA 꼬리로의 \emph{적용}이 본 작업의 기여다. 이 두 항목은
established 가 아니라 BOUNDED 로 둔다.
\end{boundbox}

% ====================================================================
\section{자기참조 결합과 Volterra 적분식}\label{sec:volterra}
% ====================================================================
\textbf{왜 자기참조인가(물리).} $\xi_j$ 와 $V_n$ 은 전하 보존(식~\eqref{eq:charge_balance})으로 얽혀 있다 — 진행
$\Theta$ 가 늘면 $V_n$ 이 이동하고(식~\eqref{eq:charge_balance} 미분), $V_n$ 이 이동하면 평형 목표
$\Theta_\eq(V_n)$ 가 바뀌며, 바뀐 목표가 다시 $\Theta$ 의 완화 방향을 정한다. 이 되먹임 때문에 $\Theta$ 는 \emph{자기
자신을 목표로 갖는} 적분식의 해가 된다.

\textbf{적분형 유도(single-mode 에서, 무비약).} 한 mode 의 $q$-좌표 완화는 $\dd\Theta/\dd q=(1/L)(\Theta_\eq-\Theta)$
다(식~\eqref{eq:single_kernel} 와 같은 1차 완화, rate $1/L$). 이 1계 선형 ODE 를 상수변화법(적분인자 $e^{q/L}$)으로
풀면 $\dd(\Theta e^{q/L})/\dd q=(1/L)\Theta_\eq e^{q/L}$ 이고 $0$ 에서 $q$ 까지 적분해
\begin{equation}
\Theta(q)=\underbrace{\Theta(0)\,e^{-q/L}}_{\text{초기값 자유 감쇠}}+\int_0^q \frac1L\,e^{-(q-q')/L}\,\Theta_\eq(q')\,\dd q'.
\label{eq:singlemode_integral}
\end{equation}
즉 \textbf{kernel $(1/L)e^{-(q-q')/L}$ 이 \emph{목표} $\Theta_\eq$ 에 곱해진} 형이다(상수 목표 $\Theta_\eq=\Theta_0$,
$\Theta(0)=0$ 이면 $\Theta=\Theta_0(1-e^{-q/L})$ 로 복귀 — 무결성). 이제 두 곳을 일반화한다: \textbf{(i) 자기참조} —
목표가 미지 진행에 의존, $\Theta_\eq\to\Theta_\eq(V_n(q',\Theta(q')))$; \textbf{(ii) spectrum} — 단일 $L$ 을
분포 $A_L$ 로 적분. 목표 $\Theta_\eq$ 는 \emph{모든 mode 공통}이라(식~\eqref{eq:xi_dist} 의 $\xi_{\eq,j}$ 가
$g$-무관) $L$-적분 밖으로 빠지고 적분은 kernel 에만 작용하므로, $(1/L)e^{-(q-q')/L}\to\mathcal K(q,q')=\int
A_L(1/L)e^{-(q-q')/L}\dd L$, 초기항 $\Theta(0)e^{-q/L}\to\Theta_\mathrm{init}=\int A_L\Theta_{\mathrm{init},0}(L)e^{-q/L}\dd L$
($\Theta_{\mathrm{init},0}(L)$=mode 별 초기 진행, 상수 목표 scalar $\Theta_0$ 과 구별). 그러면 인과 상한 $q$ 의 Volterra 2종이
\emph{유도}된다(``뜬금없이''가 아니라 식~\eqref{eq:singlemode_integral} 의 두 일반화):
\begin{equation}
\Theta(q)=\Theta_\mathrm{init}(q)+\int_0^q \mathcal K(q,q')\,\Theta_\eq\big(V_n(q',\Theta(q')),T\big)\,\dd q',
\qquad \mathcal K(q,q')=\int_{L_{\min}}^\infty A_L(L)\,\tfrac1L\,e^{-(q-q')/L}\,\dd L.
\label{eq:volterra}
\end{equation}
여기서 $\Theta_\mathrm{init}(q)=\int_{L_{\min}}^\infty A_L(L)\,\Theta_{\mathrm{init},0}(L)\,e^{-q/L}\,\dd L$ 는 peak 진입 시점의 초기 mode
진행 $\Theta_{\mathrm{init},0}(L)$ 이 외부 구동 없이 자유 감쇠하는 동차항이고, 미지 $\Theta$ 는 좌변과 목표
$\Theta_\eq(V_n(q',\Theta(q')))$ 를 통해서만 등장한다 — 즉 ``현재 진행이 현재 목표를 바꾸고, 그 목표가 다시
현재 진행을 정한다''는 되먹임이 자기참조의 정체다.
\begin{boundbox}
\textbf{이중 계상 경고.} kernel $\mathcal K$ 는 반드시 \emph{목표} $\Theta_\eq$ 에 곱한다. 만약 잔차
$[\Theta_\eq-\Theta]$ 에 곱하면 상수 목표 극한에서 정상상태가 $\Theta_0/2$ 로 어긋난다(완화를 두 번 세는
오류). 상수 목표 $\Theta_\eq=\Theta_0$, $\mathcal K=(1/L)e^{-(q-q')/L}$, $\Theta_\mathrm{init}=0$ 극한에서
식~\eqref{eq:volterra} 는 $\Theta(q)=\Theta_0(1-e^{-q/L})$ 로 single-mode ODE(식~\eqref{eq:single_kernel})
와 정확히 일치하며 $q\to\infty$ 서 목표 $\Theta_0$ 로 수렴한다(single-mode 무결성 점검). \emph{분포} $A_L$ 의 경우
같은 극한의 $q\to\infty$ 값은 $\big(\int_{L_{\min}}^\infty A_L\dd L\big)\Theta_0$ 이므로, $A_0\equiv1$(규격 보존, \S\ref{sec:spectrum})
일 때만 $\Theta_0$ 로 수렴한다 — $A_0\ne1$(진폭 spectrum)이면 총 scale 은 $\Theta\cdot Q_p$ 정의(\S\ref{sec:fiteq})의
용량 규격에 흡수되어 절대 scale 이 재조정된다(위 single-mode 점검은 $\delta$-collapse 한정).
\end{boundbox}
% ====================================================================
\section{자기참조식을 닫기 (1): Plan B — 직접 g-grid (항상 보존)}\label{sec:planB}
% ====================================================================
식~\eqref{eq:volterra} 에서 \emph{값}을 얻으려면 미지 $\Theta$ 를 닫아야 한다. 두 방법을 둔다.
\textbf{Plan B}(본 절)는 항상 보존되는 수치 core/validator 이고, \textbf{Plan A}(\S\ref{sec:planA})는
검증 통과 시 같은 결론을 닫힌형으로 표현하는 보조식이다. 핵심 물리 결론은 Plan B에서 먼저 확인한다.

\textbf{Plan B — 직접 $g$-grid 수치 적분.} 배리어를 격자 $\{g_m\}$ 로 이산화해
식~\eqref{eq:xi_dist} 를 시간 적분하며, 매 시점 전하 보존(식~\eqref{eq:charge_balance})으로 $V_n$ 을 풀어 결합한다:
\begin{equation*}
\xi_j(t)=\sum_m w_m\,\rho_G(g_m)\,\xi_j(g_m,t),\qquad \sum_m w_m\,\rho_G(g_m)=\int_0^\infty\rho_G\,\dd g=1.
\end{equation*}
즉 $w_m$ 은 격자 $\{g_m\}$ 의 quadrature 가중으로 식~\eqref{eq:xi_dist} 의 \emph{적분을 그대로 이산화}한 것이지
새 가중을 넣은 게 아니며($w_m\rho_G(g_m)$ 이 정규화 측도를 이뤄 합이 1), 비균일$\cdot$large-$L$ 편중 격자
(부록 \S\ref{sec:ch1_numeric})에서도 이 정규화는 재계산해 보존한다. 주어진 $\rho_G$ 와 완화 closure 하에서
이산화 오차만 남는 수치 기준해이며, Plan A 닫힌형을 이 기준해와 상대오차
$\varepsilon=\lVert\Theta_\mathrm A-\Theta_\mathrm B\rVert/\lVert\Theta_\mathrm B\rVert$ 로 대조하는 \emph{validator}
를 겸한다(식~\eqref{eq:closed}, 수치 스킴$\cdot$수렴 판정 상세는 부록 \S\ref{sec:ch1_numeric}). 저온$\to$큰
$L\to$긴 꼬리 같은 \emph{물리 결론은 (위 완화 closure 전제 하에) 닫힌형 없이도 성립}한다.

% ====================================================================
\section{자기참조식을 닫기 (2): Plan A — 선형화 $\to$ 닫힌형}\label{sec:planA}
% ====================================================================
Plan B 가 항상 보존되는 수치해라면, Plan A 는 \emph{데이터 피팅에 직접 쓸} 닫힌 식이다(검증 통과 시). \textbf{baseline
$\Theta^{(0)}$ 정의(먼저 명시)}: 자기참조를 끈($V_n$ 고정, 아래 $b=0$) 0차 해를 $\Theta^{(0)}(q)$ 로 둔다 —
single-mode 극한에선 $\Theta^{(0)}(q)=\Theta_0(1-e^{-(q-q_a)/L})$(식~\eqref{eq:singlemode_integral} 의 상수목표
해), 일반적으로는 Volterra(식~\eqref{eq:volterra})를 $V_n$-고정으로 평가한 해다. Plan A 는 이 $\Theta^{(0)}$ 둘레의
self-consistency 보정이다.

\emph{단계 0(선형화).} 식~\eqref{eq:volterra} 의 비선형성은 오직 $\Theta_\eq(V_n(q',\Theta(q')))$ 한 곳이다.
$\Theta^{(0)}$ 둘레에서 1차 전개하면 $\Theta_\eq(V_n(\Theta))\approx a(q')+b(q')\,\Theta(q')$, 표준 1차 Taylor
분해상 절편은 $a(q')=\Theta_\eq\big(V_n(\Theta^{(0)}(q'))\big)-b(q')\,\Theta^{(0)}(q')$(baseline 에서 평가한 목표값에서
1차 기울기$\times$baseline 을 뺀 상수항)이고, 기울기(자기참조 계수)는
\begin{equation}
b(q')=\frac{\partial\Theta_\eq}{\partial V_n}\,\frac{\partial V_n}{\partial\Theta}
=-\frac{Q_p}{C_\bg}\,\frac{\partial\Theta_\eq}{\partial V_n}\ \le\ 0,
\label{eq:selfcoupling}
\end{equation}
$\partial V_n/\partial\Theta=-Q_p/C_\bg$ 는 전하 보존(식~\eqref{eq:charge_balance})을 $Q_p\Theta\equiv\sum_j
Q_{j,\tot}\xi_j$($Q_p=\sum_j Q_{j,\tot}$=총 phase 용량), $C_\bg\equiv\partial Q_\bg/\partial V_n$(배경 chemical
capacitance, 상세 \S\ref{sec:fiteq})로 묶은 $Q_\cell q=Q_\bg(V_n)+Q_p\Theta$ 를 $\Theta$ 로 미분해 얻고($C_\bg>0$), $\partial\Theta_\eq/\partial V_n>0$(평형 진행률은 전위에 증가, 식~\eqref{eq:dxidV} 의 $\dd\xi_\eq/\dd V_n>0$
를 $\Theta_\eq=Q_p^{-1}\sum_j Q_{j,\tot}\xi_{\eq,j}$ 로 합산)이므로 $b\le0$ 이 \emph{확정}된다. $b\le0$ 의 \emph{의미}는
``진행이 늘면 목표가 줄어 진행을 \emph{되끌어당긴다}''는 음의 되먹임이다.
이로써 식~\eqref{eq:volterra} 가 미지 $\Theta$ 에 \emph{선형}이 된다: $\Theta=c+\int_0^q\mathcal K\,b\,\Theta\,\dd q'$,
$c(q)=\Theta_\mathrm{init}+\int_0^q\mathcal K\,a\,\dd q'$.
\emph{단계 1(ratio ansatz).} 사용자 PhD \cite{lee2011,son2013} 의 ratio-substitution 을 \emph{미지} $\Theta(q')$
에만 적용한다: $\Theta(q')\approx\Theta(q)\,\Theta^{(0)}(q')/\Theta^{(0)}(q)$ (진행의 \emph{모양}은 baseline 을
따르고 크기만 미지로 둔다). \emph{정당화}: 자기참조항 $\int\mathcal K b\,\Theta$ 가 $c$ 에 대해 섭동($|b|$ 작음)일
때, 형상은 0차 baseline $\Theta^{(0)}$ 이 지배하고 $b$ 는 진폭만 재조정한다 — $b\to0$ 에서 정확(식~\eqref{eq:closed}
이 $c$ 로 복귀), $|b|$ 가 커질수록 형상 왜곡이 누적되어 아래 boundbox 의 부정확 한계로 이어진다.
\emph{단계 2($\Theta(q)$ 인수분해).} 단계 1 을 단계 0 의 선형식 $\Theta=c+\int_0^q\mathcal K b\,\Theta\,\dd q'$ 의
적분에 넣으면, 앞 인자 $\Theta(q)/\Theta^{(0)}(q)$ 가 $q'$ 에 무관이라 적분 \emph{밖}으로 빠진다:
\begin{equation*}
\int_0^q\mathcal K(q,q')\,b(q')\,\Theta(q')\,\dd q'\approx\frac{\Theta(q)}{\Theta^{(0)}(q)}\int_0^q\mathcal K(q,q')\,b(q')\,\Theta^{(0)}(q')\,\dd q'.
\end{equation*}
이를 $\Theta(q)=c(q)+(\text{위 적분})$ 에 넣고 $\Theta(q)$ 항을 좌변으로 모아 분모로 묶으면:
\begin{equation}
\boxed{\;\Theta_\mathrm A(q)=\frac{c(q)}{\,1-\dfrac{1}{\Theta^{(0)}(q)}\displaystyle\int_0^q\mathcal K(q,q')\,b(q')\,\Theta^{(0)}(q')\,\dd q'\,}\;.}
\label{eq:closed}
\end{equation}
\textbf{이것이 피팅에 쓰는 닫힌형이다} — 각 항의 의미: 분자 $c$ = 자유 감쇠($\Theta_\mathrm{init}$) + baseline
목표 기여, 분모의 적분 = 자기참조 감쇠($b<0$ 이라 분모 $>1$ → 꼬리를 줄임). \emph{무결성 점검}: 자기참조가
없으면($b\to0$) 분모 $\to1$, $\Theta_\mathrm A\to c$ 로 식~\eqref{eq:volterra} 의 올바른 극한(상수 목표서
$\Theta_0(1-e^{-q/L})$, 식~\eqref{eq:single_kernel})으로 복귀한다.
\begin{boundbox}
\textbf{Fredholm$\leftrightarrow$Volterra 가정차(\AL{16}; 그대로 쓰면 안 됨).} 사용자 논문(\emph{JCP}
\textbf{147}, 144111 (2017), Eq.(32)$\to$(34), \cite{jcp2017})은 \emph{Fredholm} 2종(고정 구간, 정상상태 확률)을
ratio-substitution 으로 닫는다. graphite 의 식~\eqref{eq:volterra} 는 \emph{Volterra}(인과 상한 $q$)다. 따라서
논문 결과를 그대로 옮기지 않고, 인과성을 유지한 채 자기참조만 선형화(단계 0)해 \emph{같은} ratio 구조를 적용했다.
또 \textbf{논문 자신이 ``reaction zone 이 넓으면 ratio 근사가 부정확''이라 경고}하는데, graphite 에서 넓은
spectrum $=$ 저온 stretched tail(가장 관심 영역)이므로 closure 가 가장 필요한 곳에서 가장 부정확할 수 있다.
$\Rightarrow$ \textbf{Plan A 단독 채택 금지}: Plan B(g-grid) 대비 상대오차
$\varepsilon=\lVert\Theta_\mathrm A-\Theta_\mathrm B\rVert/\lVert\Theta_\mathrm B\rVert$ 가 운용 $T\cdot$rate 에서
작을 때만 닫힌형으로 승격한다(BOUNDED).
\end{boundbox}

% ====================================================================
\section{ICA 와 DVA, 그리고 피팅 가능한 닫힌 논리식}\label{sec:fiteq}
% ====================================================================
\textbf{내부 전위 미분.} 식~\eqref{eq:charge_balance} 를 $q$ 로 미분하면(등온)
$Q_\cell=\dfrac{\partial Q_\bg}{\partial V_n}\dfrac{\dd V_n}{\dd q}+\sum_j Q_{j,\tot}\dfrac{\dd\xi_j}{\dd q}$,
따라서
\begin{equation}
\frac{\dd V_n}{\dd q}=\frac{Q_\cell-\sum_j Q_{j,\tot}\,\dd\xi_j/\dd q}{\partial Q_\bg/\partial V_n}.
\label{eq:dVdq}
\end{equation}
전체 phase progress 를 $Q_p\Theta\equiv\sum_j Q_{j,\tot}\xi_j$ 로 묶는다($Q_p=\sum_j Q_{j,\tot}$=용량 scale,
$\Theta\in[0,1]$ 무차원). 이 $\Theta$ 는 \S\ref{sec:kernel} 의 $\Theta$ 와 \emph{같은 양}이며, 그 post-peak
도함수가 $\dd\Theta_\mathrm{tail}/\dd q$(식~\eqref{eq:kernel_integral})다. 외부 전하 $\dd Q=\dd Q_\ext=
Q_\cell\,\dd q$ 는 배경 저장 $\dd Q_\bg=C_\bg\,\dd V_n$($C_\bg\equiv\partial Q_\bg/\partial V_n$)과
phase progress $Q_p\,\dd\Theta$ 의 합이다:
\begin{equation}
\dd Q=C_\bg\,\dd V_n+Q_p\,\dd\Theta.
\label{eq:dQsplit}
\end{equation}
이는 식~\eqref{eq:dVdq} 양변에 $C_\bg=\partial Q_\bg/\partial V_n$ 을 곱하고 $\dd q=\dd Q/Q_\cell$ 로 다시 쓴
\emph{동일식}이다(새 가정이 아니라 같은 전하보존 미분의 다른 표현). 양변을 $\dd Q$ 로 나누면 $1=C_\bg\,(\dd V_n/\dd Q)+Q_p\,(\dd\Theta/\dd Q)$ 이고, $\dd V_n/\dd Q$ 에 대해 풀어
역수를 취하면 내부 전위 기준 ICA
\begin{equation}
\boxed{\;\frac{\dd Q}{\dd V_n}=\frac{C_\bg(V_n,T)}{\,1-Q_p\,(\dd\Theta/\dd Q)\,}\;.}
\label{eq:fiteq}
\end{equation}
식~\eqref{eq:kernel_integral} 의 꼬리 기여가 $\dd\Theta/\dd Q$ 에 들어가 ICA 꼬리로 나타난다(좌표 bridge:
외부 전하 $Q=Q_\ext=Q_\cell q$ 이므로 $\dd\Theta/\dd Q=(1/Q_\cell)\,\dd\Theta/\dd q$). \emph{적용 영역}: post-peak
구간에선 꼬리 기여가 지배해 $\dd\Theta/\dd q\simeq\dd\Theta_\mathrm{tail}/\dd q$ 이고, peak 영역에선
$\dd\Theta/\dd q\to\dd\Theta_\eq/\dd q$(빠른 kinetics 극한, \S\ref{sec:kernel} 직전 single-mode 논의)로 환원된다.
$\dd\Theta/\dd Q$ 에 넣는 $\Theta$ 는 닫힌 $\Theta_\mathrm A$(식~\eqref{eq:closed}) 또는 Plan B 수치해를 $Q$ 로
미분한 값이다. 분모
$1-Q_p\,(\dd\Theta/\dd Q)>0$ 은 단조성(식~\eqref{eq:monotone}, $C_\bg=\partial Q_\bg/\partial V_n>0$)과
\emph{동치가 아니라 별개}의 유효범위 조건이다: 전자는 배경 용량의 국소 단조성(정적 chemical capacitance
양수)이고, 분모 양수성은 식~\eqref{eq:fiteq} 의 역수 $\dd V_n/\dd Q=(1-Q_p\,\dd\Theta/\dd Q)/C_\bg$ 가
양이라는 \emph{동역학 경로 조건}($C_\bg>0$ 전제하 $\dd V_n/\dd Q>0$)이다. 두 조건이 함께 성립해야
$\dd Q/\dd V_n>0$ 이다. staging peak 에서 분모가
$0$ 에 접근하면 $\dd Q/\dd V_n$ 이 발산하는데 이것이 정상적인 ICA peak 이며, 그 너머(분모 $<0$)는 전압 해의
특이점으로 보아 피팅에서 제외한다. 측정 축으로는
apparent 전위(식~\eqref{eq:Vapp})로 환산한다: 등온 정전류에서 $\dd V_{n,\app}/\dd q=\dd V_n/\dd q+
s_I|I|\,\partial R_n/\partial q$, 그리고 $\dd Q_\ext/\dd V_{n,\app}=Q_\cell/(\dd V_{n,\app}/\dd q)$,
$\dd V_{n,\app}/\dd Q_\ext$ 가 그 역수다(DVA). \emph{전압축 특성 꼬리 길이}는 다음과 같이 환산된다: post-peak 에서
$\dd V_n/\dd q$ 를 $q_a$ 둘레 상수로 근사하면(식~\eqref{eq:single_kernel} 유도의 가정 — 꼬리에서
$\dd\mathcal A_j/\dd q\simeq0$ 인 완만 영역) $V_n-V_n(q_a)\simeq(\dd V_n/\dd q)|_{q_a}(q-q_a)$ 이고,
single-mode 꼬리의 지수 $e^{-(q-q_a)/L}$ 가 $\exp[-(V_n-V_n(q_a))/L_\varphi]$ 로 옮겨져
\begin{equation}
L_\varphi\equiv\big|\dd V_n/\dd q\big|_{q_a}\,L
\label{eq:Lphi}
\end{equation}
가 정의된다($\dd V_n/\dd q$ 가 꼬리에서 완만하다는 가정 위의 BOUNDED 근사; 기호표에 $L_\varphi$[V] 등재).

\textbf{피팅 평가 순서(deliverable).} 식~\eqref{eq:fiteq} 는 다음 순서로 계산되는 닫힌 논리식이다:
\begin{enumerate}[label=\textbf{S\arabic*.},leftmargin=2.4em]
\item 전하 보존(식~\eqref{eq:charge_balance})으로 $V_n(q,T,\{\xi_j\})$ 를 구한다.
\item 평형 목표 $\xi_{\eq,j}(V_n,T)$(식~\eqref{eq:logistic}).
\item 속도 $k_j=k_0 e^{-(G-\chi_j\mathcal A_j)/RT}$, $\mathcal A_j=s_{\phi,j}F(V_{n,\drive}-U_j)$
  (식~\eqref{eq:kact}; 병목 시 식~\eqref{eq:kphys}).
\item 배리어$\to$길이 $L(G)$(식~\eqref{eq:L_of_G}) $\to$ spectrum $A_L$(식~\eqref{eq:spectrum}).
\item 꼬리 중첩 $\dd\Theta_\mathrm{tail}/\dd q$(식~\eqref{eq:kernel_integral}); 자기참조는 \emph{본 장에서} 닫는다 —
  Plan B(g-grid 수치, 항상) 또는 Plan A 닫힌형(식~\eqref{eq:closed}, $\varepsilon$ 통과 시).
\item $\dd Q/\dd V_n$(식~\eqref{eq:fiteq}) $\to$ 측정축 $\dd Q/\dd V_{n,\app}$.
\end{enumerate}
\textbf{피팅 파라미터}: $\{U_j,w_j,Q_{j,\tot},Q_\bg\}$(평형, 저속 OCV/GITT), $\chi_j$(여러 $|I|$/전위),
$\{\Delta H_a,\Delta S_a,k_0\}$(Eyring-corrected 여러 $T$), $\rho_G$(독립 입력 — 역산 금지). \textbf{유효범위}:
$G-\chi_j\mathcal A_j\ge0$(Marcus, \AL{15}), $V_{n,\drive}\approx U_j+O(w_j)$. (전체 spectrum$\cdot$자기참조를
쓰는 정밀 평가의 solver 스킴$\cdot$수렴 판정은 부록 \S\ref{sec:ch1_numeric}.)

\subsection{심플 실데이터 피팅 근사식 (single-mode tail; 본 장만으로 사용 가능)}\label{sec:ch1_simplefit}
전체 spectrum 적분(식~\eqref{eq:kernel_integral}) 없이, \emph{하나의 우세 mode} 가 꼬리를 지배하는 극한
($A_L(L)=\Theta_0\delta(L-L^\ast)$; amplitude-bearing single-mode)에서는 위 모든 단계가 \textbf{종이와 연필로 닫히는} 가장 단순한 근사식으로 줄어든다. 이것이
실데이터 1차 피팅의 출발점이다. post-peak($q>q_a$) 에서 목표가 거의 정지하면(식~\eqref{eq:single_kernel}) 진행은
\begin{equation}
\Theta_\mathrm{tail}(q)\simeq\Theta_0\Big(1-e^{-(q-q_a)/L}\Big),\qquad
\frac{\dd\Theta_\mathrm{tail}}{\dd q}\simeq\frac{\Theta_0}{L}\,e^{-(q-q_a)/L},\qquad
\boxed{\;L=\frac{|I|}{Q_\cell\,k_j}=\frac{|I|}{Q_\cell k_0}\,e^{(G-\chi_j\mathcal A_j)/RT}\;}
\label{eq:simplefit}
\end{equation}
이고, 이를 fiteq(식~\eqref{eq:fiteq})에 넣으면 ICA 꼬리는 \emph{특성 길이 $L$ 의 지수 감쇠}로 나타난다.
\textbf{데이터 피팅 절차.}
\textbf{(0) 무대 고정(식별성 전제).} \S\ref{sec:falsify} 식별성 순서대로, 먼저 저속 OCV/GITT 로
$\{U_j,w_j,Q_{j,\tot},Q_\bg\}$ 를, 펄스/전류 차단으로 $R_n$ 을 고정한다($R_n$ 과 $\chi_j$ 가 공선형이므로 $R_n$ 을
\emph{먼저} 묶지 않으면 아래가 식별 불가). 측정 $V_{n,\app}$ 는 $V_n=V_{n,\app}-s_I|I|R_n$(식~\eqref{eq:Vapp} 역산,
$R_n$ 기지)으로 내부 $V_n$ 으로 \emph{직접} 환산한다(self-consistent 루프 아님).
\textbf{(1) $L$ 추출.} 꼬리 시작점 $q_a$ 를 \emph{먼저} 고정한다(ICA peak 위치, 또는 $\dd\xi_{\eq}/\dd q$ 가 임계
이하로 떨어지는 첫 $q$) — $q_a$ 를 자유 파라미터로 두면 $L$ 과 공선형이다. 그 뒤 post-peak ICA 꼬리를
모델 $y(q)=A\,e^{-(q-q_a)/L}+(\text{baseline})$ 로 회귀해 \emph{한 숫자} $L$ 을 뽑는다 — 진폭 $A$ 는 nuisance
($\Theta_0\cdot Q_p/C_\bg$ 등을 흡수; 관심량은 $L$), 단일-mode 꼬리에선 fiteq 분모 $1-Q_p\,\dd\Theta/\dd Q\!\approx\!1$
이라 ICA 꼬리 $\approx(Q_p/C_\bg)(\Theta_0/L)e^{-(q-q_a)/L}$ 로 닫힌다(전압축이면 $L_\varphi$, 식~\eqref{eq:Lphi}).
\textbf{(2) $\boldsymbol{\chi_j}$ 먼저 추출.} \emph{peak 근방}($\mathcal A_j=s_{\phi,j}F(V_{n,\drive}-U_j)$ 가 유한한 영역)에서 내부 전위
$V_{n,\drive}=V_n$(0단계 환산)를 가로축으로 삼아 여러 전류/전위 조건에서 얻은 $L$ 의 기울기
$\partial\ln L/\partial V_{n,\drive}=-\chi_j s_{\phi,j}F/RT$ 로부터 $\chi_j$ 를 먼저 얻는다. \textbf{비순환 보장}: $\mathcal A_j$ 는
각 데이터점에서 (0)단계 환산 $V_n$ 과 (0)단계 고정 $U_j(T)$ 로 계산되는 기지값이며, 이 단계에서는 $\Delta H_a$ 를 아직 쓰지 않는다.
따라서 $\chi_j$ 추출은 Arrhenius fit 과 독립인 shape/slope extraction 이다.
\textbf{(3) 유효 enthalpy(prefactor$\cdot U_j(T)$ 보정).} 이제 (2)에서 얻은 $\chi_j$ 를 고정하고 온도별 $L(T)$ 를 해석한다.
$1/L\propto k_j=k_0(T)e^{-(G-\chi_j\mathcal A_j)/RT}$ 이고 $k_0=(k_BT/h)\kappa$(Eyring, 식~\eqref{eq:kact})에
\emph{T-선형 prefactor} 가 있어, $\ln(1/L)$ 를 그대로 $1/T$ 로 회귀하면 prefactor 의 $\ln T$ 항이 기울기를 오염시킨다.
표준 Eyring 선형화($\ln(k/T)$ 대 $1/T$)대로 prefactor 를 보정하고, 또한 $\mathcal A_j=s_{\phi,j}F(V_{n,\drive}-U_j(T))$ 의
$U_j(T)$ 의존 때문에 $\chi_j\mathcal A_j(T)/RT$ 가 $1/T$ 에 섞이므로 이를 점별 기지값으로 빼낸다. 즉
\begin{equation*}
y(T)\equiv\ln\!\frac{1}{L\,T}-\frac{\chi_j\,\mathcal A_j(T)}{RT}\quad\text{를}\quad \frac1T\ \text{로 회귀} \;\Rightarrow\; \text{기울기}=-\frac{\Delta H_a}{R}\ (\text{순수}).
\end{equation*}
부호가 \emph{minus} 인 이유는 $\ln(1/L)$ 안에 이미 $+\chi_j\mathcal A_j/(RT)$ 가 들어 있기 때문이다. 이 항을 더하면 전위 보조 효과를 두 번 더하는 오류가 된다.
$\kappa(T)$ 의 약한 T-의존을 무시하는 근사는 $\Delta H_a/(RT)\gg1$($\Delta H_a\gtrsim40$ kJ/mol)에서 BOUNDED 다.
(보정 전 기울기 $-(\Delta H_a-\chi_j\mathcal A_j)/R$ 는 \emph{유효} enthalpy 로만 해석한다.)
\textbf{즉 $\{L,\chi_j,\Delta H_a\}$ 가 단일 지수 꼬리 피팅 + 식별성 순서로 나온다}(필요한 근거가 모두 본 장 안에 있다:
근사식~\eqref{eq:simplefit}, 식별성 \S\ref{sec:falsify}). 꼬리가 단일 지수에서 벗어나(stretched) 휘면 그때 비로소 spectrum(식~\eqref{eq:spectrum})$\cdot$자기참조(식~\eqref{eq:closed})가
필요하다 — \emph{단일 지수 근사의 실패 자체가 분포$\cdot$결합의 신호}다(반증, \S\ref{sec:falsify}).
\textbf{C-rate 기준식(0.2C anchor).} 전류 의존 분극을 저율 기준 곡선에 상대화하면, $0.2$C($I_{0.2C}>0$)를 anchor 로
\begin{equation}
V_{n,\app}(q;|I|)\approx V_{n,0.2C}(q)+s_I\big[\,|I|\,R_n(q,T,|I|)-I_{0.2C}R_n(q,T,I_{0.2C})\,\big],
\label{eq:02c}
\end{equation}
전류 의존이 약하면 단순형 $V_{n,\app}\approx V_{n,0.2C}+s_I(|I|-I_{0.2C})R_n(q,T)$ 다 — $0.2$C 에서 상변이 지연이
충분히 작거나 그 지연을 기준 곡선에 흡수해도 되는 경우의 실용 근사. 여기에 식~\eqref{eq:simplefit} 의 $L(|I|)$ 해석을 더할 때는 단일 방향으로 단정하지 않는다.
$k_j$가 고정이면 $L=|I|/(Q_\cell k_j)$라서 $q$-axis tail은 길어질 수 있고, 반대로 current-induced potential shift가 $k_j$를 충분히 키우면 $L$은 짧아질 수 있다.
Measured voltage axis에서는 apparent polarization이 peak 위치와 overlap도 바꾼다. 따라서 C-rate trend는 current prefactor, potential-assisted rate increase, apparent-axis polarization의 competition으로 정량 설명한다.
\emph{피팅 절차와의 연결}: 식~\eqref{eq:02c} 는 여러 $|I|$ 의 $V_{n,\drive}$ 를 공통 저율 기준으로
정렬해 위 (2)단계 $\chi_j$ 회귀의 가로축 입력을 일관되게 만드는 데 쓴다.

\textbf{초기값(EMG 등 경험식).} 비선형 피팅의 \emph{초기값}으로 $\dd Q/\dd V=B(V)+\sum_j A_j\,g_\mathrm{EMG}
(V;\mu_j,\sigma_j,\lambda_j)$ 같은 peak 함수를 쓸 수 있으나 주 모델이 아니다 — 꼬리는 $k_j,\rho_G,R_n\cdot$전하
보존의 결합 결과이며, EMG 의 $\lambda_j$(지수 꼬리율)는 식~\eqref{eq:simplefit} 의 $1/L$ 의 경험적 대응물로만 본다.

% ====================================================================
\section{온도 의존 꼬리의 분해와 반증(falsification)}\label{sec:falsify}
% ====================================================================
온도에 따른 꼬리 증감을 한 항으로 단정하지 않고 \textbf{세 계층}의 결합으로 본다.
\begin{longtable}{p{0.24\textwidth} p{0.30\textwidth} p{0.36\textwidth}}
\toprule 계층 & 담당 항 & 해석 \\ \midrule \endhead
평형 열역학 & $U_j(T),w_j(T),Q_\bg(V_n,T)$ & 전이 중심 전위$\cdot$평형 폭$\cdot$배경 용량의 온도 의존 \\
동역학 지연 & $k_j(T,I),\rho_G(g;T)$ & 진행률이 평형을 못 따라가 생기는 꼬리/broadening \\
전압축 분극 & $R_n(q,T,|I|)$ & apparent 전위 이동, 기울기 변화, 분극 peak 변형 \\
\bottomrule
\end{longtable}
\textbf{전위 의존 spectrum shift.} 식~\eqref{eq:L_of_G} 에서 $\partial\ln L/\partial V_{n,\drive}$
$=-\chi_j s_{\phi,j}F/RT<0$(방전 방향, $\chi_j\ge0$): 구동 전위만 독립적으로 커지면 $L$ 이 작아져 꼬리가 짧아진다.
C-rate trend는 앞 절처럼 current prefactor와 apparent-axis polarization까지 함께 보아야 한다.

\textbf{반증 규칙(forward-prediction 검정; 역산 금지).} 본 이론은 다음으로 \emph{반증 가능}하다.
\begin{enumerate}[label=(N\arabic*),leftmargin=2.4em]
\item \textbf{평형 형태}: 저속 OCV 의 near-peak 감쇠가 $\log(dQ/dV)$ 대 $(V-U)$ 면 logistic, 대 $(V-U)^2$
  면 가우시안 — 어느 쪽도 안 맞으면 평형 ansatz 재검토.
\item \textbf{Arrhenius}: Eyring prefactor와 $\mathcal A_j(T)$를 보정한 $y(T)=\ln[1/(LT)]-\chi_j\mathcal A_j/(RT)$가 $1/T$에 대해 직선이 아니면 단일 활성화 그림 기각.
\item \textbf{전위 의존}: $\partial\ln L/\partial V_{n,\drive}$ 가 위 예측과 다르면 $\chi_j$ 해석 기각.
\item \textbf{비유일성 차단}: $\rho_G$ 를 꼬리에서 역산하지 않고 독립 입력으로만 — 역산하면 비유일이라
  반증력이 없다.
\item \textbf{비퇴화 — transport$\cdot$charge-transfer 와의 분리(★ 본 이론의 가장 약한 지점).} 꼬리 길이 scaling
  $L\propto|I|$ 과 ``전위 커지면 꼬리 짧아짐''은 \emph{단순 transport 분극$\cdot$charge-transfer 도 공유}하므로,
  그것만으로 barrier-spectrum 에 귀속하면 \emph{퇴화}(오귀속)다. 구별 signature(두 조건 (a)$\cdot$(b)가 \emph{모두}
  필요하며, $k_{j,\lim}$ 이 전위무관일 때만 성립): (a) \emph{전위 의존 Arrhenius slope} — 보정 전 측정 기울기
  $\Delta H_a^{\eff}=\Delta H_a-\chi_j\mathcal A_j$(식~\eqref{eq:simplefit}; 추출은 \S\ref{sec:ch1_simplefit} 절차 (2))
  가 구동 전위에 따라 이동함. \emph{단 이것만으로는 transport 와 여전히 퇴화할 수 있다}: 고체상 확산이 율속이고
  확산계수가 점유율 의존($D=D(\xi)$, thermodynamic factor 경유)이면 transport-limited 유효 엔탈피도 전위에 따라
  이동한다(식~\eqref{eq:kphys} 의 $k_{j,\lim}$ 전위의존). 따라서 (a)는 GITT relaxation transient 의 형상
  ($\sqrt t$ Cottrell=확산 율속 vs 지수=계면 율속)으로 $k_{j,\lim}$ 의 전위의존을 독립 판정해 \emph{배제}한 뒤에만
  barrier-lowering signature 로 유효하다; (b) charge-transfer 과전압 $\eta_{\mathrm{ct}}$ 도 전위 지수의존이라 단일
  전극서 $\chi_j$ 와 공선형이므로 — \textbf{GITT/전류 차단으로 $\eta_{\mathrm{ct}}(R_{\mathrm{ct}})$ 를 독립 분리한
  뒤에만} $\chi_j$ 의 ``배리어 낮춤'' 귀속이 성립한다(단순 rate-broadening \cite{dubarry2022,fly2020} 과의 구별이
  핵심). 이 분리(확산 율속 배제 + $\eta_{\mathrm{ct}}$ 분리) 없이 단일 데이터셋에서 꼬리를 barrier-spectrum 으로
  귀속하는 것은 금지한다.
\end{enumerate}
\textbf{파라미터 제한 순서(식별성).} 모든 파라미터를 한 데이터에서 동시에 자유 피팅하면 공선형(co-linear)이라
분리 불가다. 다음 \emph{순차} 제약으로 식별성을 확보한다: \textbf{(1)} 저속 OCV/GITT(준평형, $|I|\to0$)로
$\{U_j,w_j,Q_{j,\tot},Q_\bg\}$ 를 먼저 고정 — 동역학 꼬리가 사라진 평형 무대를 잡는다. \textbf{(2)} 펄스/전류
차단으로 $R_n$(전압축 분극)을 독립 제한 — $R_n$ 과 $\chi_j$ 는 둘 다 꼬리를 늘려 공선형이므로 $R_n$ 을 먼저
묶는다. \textbf{(3)} 여러 전류/전위($\partial\ln L/\partial V_{n,\drive}$)로 $\chi_j$ 를 먼저 추출한다.
\textbf{(4)} 그 $\chi_j$ 로 $y(T)=\ln[1/(LT)]-\chi_j\mathcal A_j/(RT)$ 를 만든 뒤 여러 온도의 Eyring-corrected Arrhenius($y(T)$ 대 $1/T$)로
$\{\Delta H_a,\Delta S_a,k_0\}$ 를 제한한다. \textbf{(5)} $\rho_G$ 는 독립 입력(역산 금지, N4).
이 순서를 어기면 식별성이 무너져 \emph{어떤 적합도라도} 물리 해석이 불가능하다(과적합).

% ====================================================================
\section{부록: 자기참조식의 수치 평가와 검증 (Plan B 상세)}\label{sec:ch1_numeric}
% ====================================================================
전체 spectrum$\cdot$자기참조를 쓰는 정밀 평가는 Plan B(직접 $g$-grid)로 푼다. 단순 꼬리는 본문
\S\ref{sec:ch1_simplefit} 의 단일 지수 근사로 충분하나, stretched 꼬리(저온)나 강한 자기참조에서는 본 수치
평가가 기준해다.
\textbf{스킴.} (i) 배리어 격자 $\{g_m\}_{m=1}^{N_g}$ 를 관심 $T\cdot$전위에서 $L(g_m)$(식~\eqref{eq:L_of_G})이
관측 윈도를 덮도록 — 특히 large-$L$(긴 꼬리) 영역 — 잡는다. (ii) mode 별 ODE
$\dd\xi_j(g_m)/\dd t=k_j(g_m)[\xi_{\eq,j}-\xi_j(g_m)]$ 를 시간 적분하되, 매 시점 전하 보존
(식~\eqref{eq:charge_balance})의 대수 제약으로 $V_n$ 을 root-find 한다 — 미분변수 $\xi_j(g_m)$ + 대수 제약
$V_n$ 의 \emph{index-1 DAE} 이며, 제약을 한 번 미분하면 $\partial Q_\bg/\partial V_n\ge\epsilon_Q>0$
(식~\eqref{eq:monotone})이 $\dd V_n/\dd t$ 를 유일하게 주므로 index 가 정확히 1 이다. (iii)
$\xi_j(t)=\sum_m w_m\rho_G(g_m)\xi_j(g_m,t)$ 로 평균.
\textbf{근거 기반 tolerance(임의 상수 아님).} root-find 종료는 전하 보존 잔차 $<\delta_q Q_\cell$($\delta_q$=좌표
분해능), 시간적분은 $\xi\in[0,1]$ 상대오차와 측정 ICA 잡음 floor 에서 잡는다 — cap/clip 이 아닌 측정$\cdot$물리
scale 기준.
\textbf{Plan A 승격 검증.} 닫힌형 $\Theta_\mathrm A$(식~\eqref{eq:closed})는 Plan B 대비
$\varepsilon=\lVert\Theta_\mathrm A-\Theta_\mathrm B\rVert/\lVert\Theta_\mathrm B\rVert$ 를 운용 $T\cdot$rate$\cdot$
tail-window 에서 측정해 $\varepsilon\le\varepsilon_{\mathrm{tol}}=\max(\varepsilon_{\mathrm{disc}},\varepsilon_{\mathrm{noise}})$
(기준해 이산화 오차 $\vee$ 측정 잡음)일 때만 승격한다. \emph{특정 수치값}은 데이터 의존이라 미리 고정하지 않는다(FLAGGED). 넓은
spectrum(저온 stretched)에서 ratio 근사가 가장 열화하므로(\cite{jcp2017} 자기명시) 이 검증이 필수다.
\textbf{격자 수렴.} $\lVert\Theta_{N_g}-\Theta_{2N_g}\rVert/\lVert\Theta_{2N_g}\rVert\le\varepsilon_{\mathrm{disc}}$
(Plan A 승격에 쓰는 이산화 오차 기준 $\varepsilon_{\mathrm{disc}}$ 와 동일값)가 될 때까지 $N_g$ 를 배가한다 —
정성적 ``수렴할 때까지''가 아니라 위 $\varepsilon_{\mathrm{tol}}$ 체인과 닫힌 정량 종료조건이다. solver 코드 구현
자체는 본 이론 범위 밖(P1; 이론식$\cdot$검증 원리만 제시).

\section*{Chapter 2--5 로의 전달}
본 장이 확정한 상태변수($V_n,\xi_j,k_j,\xi_{\eq,j}$)와 식 위에서: Ch.2 는 \emph{전지 가역 반응열}을 평형
전위 온도계수로 잇고, Ch.3 은 $k_j$ 를 forward/backward 반응속도론(Level B)으로 확대하며, Ch.4 는 그 위에서
발열을, Ch.5 는 충방전 히스테리시스를 다룬다. \emph{(기존 Ch.6 의 피팅 실무$\cdot$수치 내용은 본 장
\S\ref{sec:ch1_simplefit}(심플 근사식)$\cdot$\S\ref{sec:planB}$\cdot$\S\ref{sec:planA}(Plan B/A closure)$\cdot$
\S\ref{sec:ch1_numeric}(수치 평가)로 흡수했다 — Ch1 자기완결.)}

\begin{thebibliography}{99}
\bibitem{doyle1993} M. Doyle, T. F. Fuller, J. Newman, ``Modeling of Galvanostatic Charge and Discharge
  of the Lithium/Polymer/Insertion Cell,'' \emph{J. Electrochem. Soc.} \textbf{140}, 1526 (1993).
  DOI: 10.1149/1.2221597.
\bibitem{newman} J. Newman, K. E. Thomas-Alyea, \emph{Electrochemical Systems}, 3rd ed., Wiley (2004).
  ISBN 978-0-471-47756-3.
\bibitem{mckinnon1983} W. R. McKinnon, R. R. Haering, ``Physical Mechanisms of Intercalation,'' in
  \emph{Modern Aspects of Electrochemistry} No.~15, Plenum, New York (1983), p.~235.
\bibitem{bazant2013} M. Z. Bazant, ``Theory of Chemical Kinetics and Charge Transfer Based on
  Nonequilibrium Thermodynamics,'' \emph{Acc. Chem. Res.} \textbf{46}, 1144 (2013).
  DOI: 10.1021/ar300145c.
\bibitem{eyring1935} H. Eyring, ``The Activated Complex in Chemical Reactions,'' \emph{J. Chem. Phys.}
  \textbf{3}, 107 (1935). DOI: 10.1063/1.1749604.
\bibitem{evanspolanyi1938} M. G. Evans, M. Polanyi, ``Inertia and Driving Force of Chemical Reactions,''
  \emph{Trans. Faraday Soc.} \textbf{34}, 11 (1938). DOI: 10.1039/TF9383400011.
\bibitem{bardfaulkner} A. J. Bard, L. R. Faulkner, \emph{Electrochemical Methods}, 2nd ed., Wiley (2001).
  ISBN 978-0-471-04372-0.
\bibitem{marcus1956} R. A. Marcus, ``On the Theory of Oxidation-Reduction Reactions Involving Electron
  Transfer. I,'' \emph{J. Chem. Phys.} \textbf{24}, 966 (1956). DOI: 10.1063/1.1742723.
\bibitem{onsager1931} L. Onsager, ``Reciprocal Relations in Irreversible Processes. I,'' \emph{Phys. Rev.}
  \textbf{37}, 405 (1931). DOI: 10.1103/PhysRev.37.405.
\bibitem{degrootmazur1962} S. R. de Groot, P. Mazur, \emph{Non-Equilibrium Thermodynamics},
  North-Holland (1962); Dover (1984).
\bibitem{dahn1991} J. R. Dahn, ``Phase Diagram of Li$_x$C$_6$,'' \emph{Phys. Rev. B} \textbf{44}, 9170
  (1991). DOI: 10.1103/PhysRevB.44.9170.
\bibitem{ohzuku1993} T. Ohzuku, Y. Iwakoshi, K. Sawai, ``Formation of Lithium-Graphite Intercalation Compounds in Nonaqueous Electrolytes,''
  \emph{J. Electrochem. Soc.} \textbf{140}, 2490 (1993). DOI: 10.1149/1.2220849.
\bibitem{funabiki1999ea} A. Funabiki, M. Inaba, T. Abe, Z. Ogumi, ``Nucleation and Phase-Boundary
  Movement upon Stage Transformation in Lithium-Graphite Intercalation Compounds,''
  \emph{Electrochim. Acta} \textbf{45}, 865 (1999). DOI: 10.1016/S0013-4686(99)00290-X.
\bibitem{funabiki1999jes} A. Funabiki, M. Inaba, T. Abe, Z. Ogumi, ``Stage Transformation of
  Lithium-Graphite Intercalation Compounds Caused by Electrochemical Lithium Intercalation,''
  \emph{J. Electrochem. Soc.} \textbf{146}, 2443 (1999). DOI: 10.1149/1.1391953.
\bibitem{baessler1993} H. B\"assler, ``Charge Transport in Disordered Organic Photoconductors,''
  \emph{Phys. Status Solidi B} \textbf{175}, 15 (1993). DOI: 10.1002/pssb.2221750102.
\bibitem{svare2000} I. Svare, S. W. Martin, F. Borsa, ``Stretched Exponentials with $T$-Dependent
  Exponents from Fixed Distributions of Energy Barriers for Relaxation Times in Fast-Ion Conductors,''
  \emph{Phys. Rev. B} \textbf{61}, 228 (2000). DOI: 10.1103/PhysRevB.61.228.
\bibitem{johnston2006} D. C. Johnston, ``Stretched Exponential Relaxation Arising from a Continuous Sum
  of Exponential Decays,'' \emph{Phys. Rev. B} \textbf{74}, 184430 (2006). DOI: 10.1103/PhysRevB.74.184430.
\bibitem{lindsey1980} C. P. Lindsey, G. D. Patterson, ``Detailed Comparison of the Williams-Watts and
  Cole-Davidson Functions,'' \emph{J. Chem. Phys.} \textbf{73}, 3348 (1980). DOI: 10.1063/1.440530.
\bibitem{jcp2017} K. Lee, S. Lee, C. H. Choi, S. Lee, ``Effects of External Electric Field and Anisotropic Long-Range Reactivity on Charge Separation Probability,''
  \emph{J. Chem. Phys.} \textbf{147}, 144111 (2017). DOI: 10.1063/1.5000882.
\bibitem{lee2011} S. Lee, C. Y. Son, J. Sung, S.-H. Chong, ``Communication: Propagator for Diffusive Dynamics of an Interacting Molecular Pair,''
  \emph{J. Chem. Phys.} \textbf{134}, 121102 (2011). DOI: 10.1063/1.3565476.
\bibitem{son2013} C. Y. Son, J. Kim, J.-H. Kim, J. S. Kim, S. Lee, ``An Accurate Expression for the Rates of Diffusion-Influenced Bimolecular Reactions with Long-Range Reactivity,''
  \emph{J. Chem. Phys.} \textbf{138}, 164123 (2013). DOI: 10.1063/1.4802584.
\bibitem{dubarry2022} M. Dubarry, D. Anseán, ``Best practices for incremental capacity analysis,''
  \emph{Front. Energy Res.} \textbf{10}, 1023555 (2022). DOI: 10.3389/fenrg.2022.1023555.
\bibitem{fly2020} A. Fly, R. Chen, ``Rate dependency of incremental capacity analysis (dQ/dV) as a diagnostic tool for lithium-ion batteries,''
  \emph{J. Energy Storage} \textbf{29}, 101329 (2020). DOI: 10.1016/j.est.2020.101329.
\end{thebibliography}

\end{document}
